%--------------------------------------
% Document type, language, math support, layout and margins
\documentclass[article, 11pt, reqno]{amsart}
\usepackage[utf8]{inputenc}
\usepackage[english]{babel}
\usepackage{mathtools, amsmath, amsthm, amsfonts, amssymb} % Math packages
\usepackage{bm}

\usepackage[hidelinks]{hyperref} % for hyper references
\usepackage{cleveref}
\crefname{section}{§\hspace{-0.1cm}}{§§}
\Crefname{section}{§}{§§}

%--------------------------------------
% For controlling page margins
\usepackage[top=3cm, bottom=3cm, left=4cm, right=4cm]{geometry}
% For bibliography and citation style
\usepackage{csquotes}
\usepackage{natbib}
\setcitestyle{aysep={}}
\usepackage[toc,page]{appendix} % For linking to the appendix
%--------------------------------------
% text spacing, alignment, and formatting
\usepackage{setspace}
\usepackage[multiple,hang,flushmargin]{footmisc} % Remove indent from footnotes, allow multiple footnotes
\usepackage[foot]{amsaddr}
\setlength{\skip\footins}{.3cm} % Control spacing between footnotes and body text
\onehalfspacing
\frenchspacing
\raggedbottom
%--------------------------------------
% aesthetic tweaks for quoting, figures, and tables
\usepackage{upquote} % For `' quote formatting
\usepackage{graphicx} % For importing images
\usepackage{booktabs} % Improves table options & aesthetics
\usepackage{caption} % For captions
\usepackage{subcaption} % For sub-captions
\usepackage{makecell} %  Create multi-lined tabular cells
%--------------------------------------
% documents-specific packages
\usepackage{tikz} % For phase diagrams
\usetikzlibrary{arrows.meta} % Arrows library
\usepackage{pgfplots} % For native figures
\pgfplotsset{compat=1.17}
%--------------------------------------
\usepackage[footnote]{acronym} % For removing footnote indents
\usepackage{scrextend} % For removing footnote indents
\deffootnote[1em]{0em}{0em} % For removing footnote indents
{\textsuperscript{\thefootnotemark}\,\,}
\usepackage{multicol}
\usepackage{listings}
%--------------------------------------
\makeatother
\DeclarePairedDelimiter\ceil{\lceil}{\rceil}
\DeclarePairedDelimiter\floor{\lfloor}{\rfloor}
\DeclareMathOperator*{\argmax}{arg \ max}
%--------------------------------------
\renewcommand{\floatpagefraction}{.9}
\renewcommand{\topfraction}{.9}
\renewcommand{\bottomfraction}{.95}
\renewcommand{\textfraction}{.05}
%--------------------------------------
\newenvironment{myfont}{\fontfamily{qcr}\selectfont}{\par}
%--------------------------------------
% Some custom theorems
\newtheorem{theorem}{Theorem}
\newtheorem{prop}{Proposition}
\newtheorem{axiom}[theorem]{Axiom}
\newtheorem{case}[theorem]{Case}
\newtheorem{claim}[theorem]{Claim}
\newtheorem{conclusion}[theorem]{Conclusion}
\newtheorem{condition}[theorem]{Condition}
\newtheorem{conjecture}[theorem]{Conjecture}
\newtheorem{corollary}{Corollary}
\newtheorem{criterion}[theorem]{Criterion}
\newtheorem*{definition}{Definition}
\newtheorem{example}{Example}
\newtheorem{exercise}[theorem]{Exercise}
\newtheorem{lemma}{Lemma}
\newtheorem{notation}[theorem]{Notation}
\newtheorem{problem}[theorem]{Problem}
\newtheorem{proposition}[theorem]{Proposition}
\newtheorem{remark}{Remark}
\newtheorem{solution}[theorem]{Solution}
\newtheorem*{proprestate}{Proposition~1}
\newtheorem{summary}[theorem]{Summary}
%--------------------------------------
\newtheorem{innercustomthm}{Proposition}
\newenvironment{customthm}[1]
{\renewcommand\theinnercustomthm{#1}\innercustomthm}
{\endinnercustomthm}
\newtheorem{innercustomcor}{Corollary}
\newenvironment{customcor}[1]
{\renewcommand\theinnercustomcor{#1}\innercustomcor}
{\endinnercustomcor}
\newtheoremstyle{named}{}{}{\itshape}{}{\bfseries}{.}{.5em}{\thmnote{#3}#1}
\theoremstyle{named}
\newtheorem*{namedtheorem}{}
%--------------------------------------
\newcommand\independent{\protect\mathpalette{\protect\independenT}{\perp}}
\def\independenT#1#2{\mathrel{\rlap{$#1#2$}\mkern2mu{#1#2}}} % Define statistical independence symbol
%--------------------------------------
\usepackage[indent=0pt]{parskip} % for paragraph spacing
%--------------------------------------
\usepackage{comment} % for commenting long sections
%--------------------------------------
\raggedbottom % for ragged bottom to paragraphs
%----------------------------------------------------------------
\onehalfspacing  % for 1.5x line height
\frenchspacing

%----------------------------------------------------------------
\title[Funding the capability-resource gap]{Funding the capability-resource gap:\\ a model of optimal science funding}

\author{Aydin Mohseni}
\address{Department of Philosophy, Carnegie Mellon University}
\email{amohseni@andrew.cmu.edu}
\thanks{This work was supported by the John Templeton Foundation [grant number TODO].}

\author{Simon DeDeo}
\address{Department of Social and Decision Sciences, Carnegie Mellon University}

\author{Kevin J.\,S. Zollman}
\address{Department of Philosophy, Carnegie Mellon University}

\begin{document}

\begin{abstract}
The grant funder's challenge is to maximize the impact of a fixed budget under uncertainty
about the researchers who might receive it. We develop a Bayesian model of this challenge:
researchers differ in their epistemic capabilities and in resources; research output
requires both, and funding compounds capability across rounds. We show that the optimal
allocation funds the capability-resource gap: grants increase with capability and decrease
with existing resources. Track records underdetermine the gap; peer review supplies
an additional, noisy signal of capability. We find that this signal's value increases with
capability inequality and with its informativeness, and that under heavy-tailed capability
distributions even a coarse signal suffices. We further find that uniform seed grants and
lotteries, which forgo targeting, lose little output exactly where review is worth little, and that
whom to fund matters far more than when to fund.
\end{abstract}

\maketitle

\section{Introduction}\label{sec:intro}

Science funding is large and consequential. How best to allocate it remains poorly
understood. The United States performed \$937 billion of R\&D in 2023 and an estimated
\$993 billion in 2024, roughly a fifth of it funded by the federal government
\citep{NCSES2026, NSB2026}. In spite of this, basic dimensions of allocation are contested. Whether
funding should be concentrated on a few researchers or spread across many is a long-running
dispute \citep{Aagaard2020}. Whether peer review is effective is disputed: among funded NIH
grants, percentile scores predicted subsequent productivity barely better than chance
\citep{Fang2016}, while other designs find modest but real predictive power
\citep{Gallo2014, LiAgha2015, ParkLeeKim2015}. % [check all three]
And whether review should be replaced, in part, by lotteries has moved from proposal to
practice \citep{Roumbanis2019, Heyard2022}. % [check: SNSF; NZ HRC; Volkswagen]
These disputes concern a single underlying problem: what a funder with a fixed budget and
imperfect knowledge of researchers should do. This paper develops a model of that problem
and works out its logic: whom to fund, whether to spend early or wait and learn, what
information about researchers is worth, and what egalitarian alternatives lose.

Consider the challenge of the grant funder, whether a federal agency disbursing billions or
a private foundation disbursing millions: you have a fixed budget and must choose how to
allocate it across researchers and research projects to create the most positive impact.
There are two intuitive schemes. Fund the track record: give to those who have produced the
most, since past production is the best evidence of future production. Fund the
under-resourced: give to those who have the least, since a dollar goes further where dollars
are scarce. Each scheme can be intuitively compelling. Both, we will show, are suboptimal,
and for the same reason: each accounts for only one of the two quantities the impact of a
grant depends on.

We model this decision situation as follows. A funder with a fixed budget allocates
grants, over multiple rounds, to researchers who differ in two ways: epistemic capability,
what a researcher knows and can do, and resources, what a researcher has to work with.
Research output requires both and is limited by the scarcer of the two; the funder can
observe output, but not capability or resources themselves. A grant adds to a researcher's resources in
the round it is given and, because doing research builds expertise, compounds their
capability in later rounds. Each round, the funder observes the output that results,
updates its beliefs about the researchers, and allocates again.

Our results are as follows. In the ideal condition where the funder has complete
information regarding researchers' capabilities and resources, the optimal allocation funds
the capability-resource gap: each funded researcher receives a grant increasing in their
capability and decreasing in their existing resources, and the budget determines how many
researchers receive funding at all. Each intuitive scheme fails because it
accounts for only one of the two quantities. In reality the funder lacks direct access to
the gap, and track records alone underdetermine it: we show that a funder relying on the
track record performs far worse than one that estimates the gap.

Peer review supplies an additional signal of capability. Its value for a Bayesian funder
is the expected value of its reduction of the funder's uncertainty regarding the gap, and
that value is greatest where capability is unequally distributed and review informative; in
fields with heavy-tailed distributions of capability, even a coarse signal captures most of
this value. The same reasoning gives what the egalitarian alternatives lose: uniform seed
grants and lotteries lose the value of the targeting they displace, little where review is
worth little, most where review is worth most. The budget is a variable throughout,
so our results speak across funder scales: a small foundation, whose grants are modest
beside researchers' total resources, is the funder for whom targeting matters most per
dollar. Whether optimal funding should concentrate on a few researchers or spread across many
depends on the size of the funder's budget relative to researchers' total resources. We show
how strategic timing of when one provides funds matters far less than strategic targeting
of whom one provides funds: how much of the budget to spend early rather than late shifts
with the funder's budget relative to researchers' total resources, but the difference is
small relative to the difference between funding the right and the wrong researchers.

Two substantive assumptions bound these results. Grants are consumed, a grant's durable
residue (skills, publications, standing) being carried as capability; and capability and
resources are complements throughout, neither fully substituting for the other. We defend
each assumption where it is used.

We proceed as follows. \S\ref{sec:literature} situates our model in the literatures on
funding allocation, peer review, and lotteries. \S\ref{sec:model} presents our model. \S4 derives the
optimal allocation. \S5 shows why track records and small grants cannot reveal the gap.
\S6 measures the value of peer review and relates it to the empirical evidence on review.
\S7 examines the timing of funding, including the strongest argument for spending early.
\S8 measures what seed grants and lotteries lose and asks how widely funds should be
spread. \S9 discusses implications and limitations, and \S10 concludes.
% NOTE: \S3-\S10 hard-coded until those sections exist; convert to \S\ref{...} as each lands.

\section{Related literature}\label{sec:literature}

Our contribution sits at the intersection of four literatures: the empirical debate over
concentrating or dispersing funds, the evidence on the efficacy of peer review, the case
for funding lotteries, and formal models of science funding.

Whether funders should concentrate resources on a few researchers or spread them across
many is a long-running dispute, and concentration is the current practice: the top 1\% of
NIH principal investigators receive close to 10\% of research project grant funding, a
share that has risen since 1985 \citep{KatzMatter2020, Lauer2021}. Empirical studies of
funded portfolios generally find diminishing marginal returns to funding per researcher,
and reviews of this literature conclude that the evidence favors more dispersal than
current practice \citep{Mongeon2016, Aagaard2020}. Concentration is also self-reinforcing:
early grant winners accumulate advantages beyond the grant itself \citep{BolVaanRijt2018},
a dynamic that formal models show can reward questionable research practices over a career
\citep{Heesen2024}.
This dispute typically concerns the returns to funding: whether output per researcher increases
or decreases as grants concentrate. In our model, the answer also depends on the size of the
funder's budget relative to researchers' total resources: with a budget large enough to
relieve most bottlenecks, the optimal allocation itself reaches nearly every researcher, so
little is lost by spreading funds widely; with a tight budget, effectiveness requires
concentrating funds on the researchers with the largest capability-resource gaps (\S8).
The value of any targeting, in turn, depends on what the funder can know about researchers
(\S6).

The efficacy of peer review is contested. Among funded NIH grants, percentile scores
predict subsequent productivity barely better than chance \citep{Fang2016}, and reviewers
assessing the same application agree only weakly \citep{Pier2018}. Analyses with different
designs find real but modest predictive power: better scores predict more publications,
citations, and patents, including under controls for investigator history and institution
\citep{LiAgha2015}, and a funding natural experiment finds higher-scored proposals
outproduce lower-scored ones \citep{ParkLeeKim2015}. Both sides of this dispute measure
review's accuracy: how well scores predict outcomes among funded applicants. Our model
instead measures review's value: the expected value of its reduction of the funder's
uncertainty regarding the capability-resource gap. That value depends on the field; the same noisy signal can be
worth a great deal in one field and nothing in another (\S6), and no single number settles
the question.

Doubts about review have made lotteries a live proposal. Advocates argue that randomizing
among fundable applications saves review costs, reduces bias, and honestly reflects
review's limited ability to discriminate among fundable proposals
\citep{FangCasadevall2016, Avin2019, Roumbanis2019}; several funders have implemented
partial lotteries: the Health Research Council of New Zealand has allocated its Explorer
Grants by lottery since 2013 \citep{Liu2020}, the Swiss National Science Foundation
randomizes among ties \citep{Heyard2022}, and a German funding line runs a lottery before
review \citep{Luebber2025}.
The existing case measures what a lottery saves: \citet{GrossBergstrom2019} model grant
competition as a contest in which scientists invest costly effort in proposals, show that
when only a small fraction of proposals is funded this effort can rival the scientific
value of the research funded, and show that a partial lottery among proposals above a
threshold removes the link between the wasted effort and the fraction funded. What has not
been measured is what a lottery forgoes. Our model measures it: the value of targeting,
which runs from negligible, where budgets are ample,
signals uninformative, or capability evenly spread, to substantial, where capability is
heavy-tailed and budgets tight (\S8). Together the two measures let a funder weigh the
choice field by field.

Formal models of science funding have concentrated on the incentive costs of competition
and on the value of exploration. Beyond the contest model above, \citet{GrossBergstrom2025}
find, in a game-theoretic model of contests for scarce rewards, that intensifying
competition pushes scientists toward higher-risk, higher-return projects.
\citet{AvinBJPS2019}
simulates funding strategies on an epistemic landscape and finds that allocation with an
explicitly random component significantly outperforms a peer-review-like strategy on large
landscapes, because review concentrates effort on regions already known to be productive.
\citet{Sorzano2026} treat allocation as a decision problem under heavy-tailed uncertainty
about project value, pairing an empirical bibliometric signal with a biased-lottery
mechanism. In each of these models, funding raises a single quantity, a project's
chance of success or a researcher's productivity, and none treats the timing of spending as
the funder's choice.
In our model, a researcher's output increases in both capability and resources but is
limited by the scarcer of the two, and the funder observes neither directly.

Our formulation and analysis of the capability-resource gap allows us to address questions
raised by these literatures: how widely funds should be spread (\S8), what peer review is
worth (\S6), and when lotteries are defensible (\S8).
% NOTE: \S4-\S8 hard-coded until those sections exist; convert to \S\ref{...} as each lands.

\section{Our model}\label{sec:model}

A funder faces a population of $n$ researchers over $T$ funding rounds. Researcher $i$ has
epistemic capability $K_i$ and resources $R_i$. Initial capabilities and resources are
drawn from Pareto distributions with tail parameters $\alpha_K$ and $\alpha_R$; smaller
tail parameters produce more unequal fields, and a correlation parameter $\rho$ allows
capable researchers to tend to be well resourced. The funder knows these distributions; it
never observes any researcher's capability or resources directly.

In each round, researcher $i$ produces research output\footnote{We are agnostic about the
interpretation of output units: they may be read as publications, as findings, or as units
of scientific value.} at random, in discrete units, with expected output\footnote{Output in
each round is drawn from a Poisson distribution with rate $\lambda_i$, independently across
researchers and rounds given the state.}
\[
\lambda_i = A\,\Lambda(K_i, R_i),
\]
where $A$ is a productivity constant and $\Lambda$ is a production function that increases
in both arguments but is limited by the scarcer of the two. Our main analysis takes
$\Lambda$ to be the harmonic mean:
\[
\lambda_i = \frac{2A K_i R_i}{K_i + R_i}.
\]
The harmonic mean makes capability and resources complements without making either an
absolute requirement: a capable researcher with thin resources still produces, by working
more slowly or borrowing equipment, but no surplus of one input fully substitutes for a
shortage of the other. The harmonic mean sits between two boundary cases that fail in
opposite directions: Cobb-Douglas, under which a capable researcher can compensate almost
fully for vanishing resources, and Leontief, under which no compensation is possible at
all. Appendix~\ref{app:technology} repeats our analysis across the family between these
boundaries, and \S8 uses the family to ask how widely funds should be spread.

Capabilities increase through research. Between rounds,
\[
K_i \leftarrow K_i + \epsilon\,\Lambda(K_i, R_i),
\]
where $\epsilon$ is the compounding rate: doing research builds expertise, and the same
bottleneck that limits output limits growth. Growth requires both inputs, and capability
never declines. Resources, by contrast, do not accumulate: resources in a round are
baseline plus grant, $R_i = R_{i0} + g_i$. The baseline $R_{i0}$ recurs every round,
representing the time and support a researcher has independent of the funder, while grants
are consumed in the round they are given. This is a substantive assumption, not
bookkeeping: whatever endures from a grant, the skills, publications, and standing it
produces, is carried by the capability update, so funding's lasting effect flows through
what researchers become rather than what they keep. Science is in this respect an
accumulation process: success begets the means of further success.

The funder's total budget is fixed before any outcomes are observed. We parameterize it by
a budget scale $b$: at $b = 0.5$, the purse equals the community's expected baseline
resources for one round, and the purse scales linearly in $b$. (The exact normalization is
in Appendix [X].) % TODO: simulation-specification appendix (owed; carries the normalization, parameter table, code pointer) The budget does not grow with the horizon: a longer horizon varies the
timing of spending, never its amount. Strategies that do not plan ahead spend an equal
share of the purse each round. Forward-looking strategies may spend on any schedule: each
round, they re-plan the entire remaining budget over the remaining horizon and execute the
current round's allocation.

The funder learns as it goes. It observes each researcher's track record, the output they
have produced to date, as it accumulates, and two further signals are available. A resource
signal, observed once at the start, gives a noisy reading of each researcher's baseline
resources (noise $\tau_R$), as one might obtain from institutional information. A review
signal, drawn afresh in each round that a strategy uses it, gives a noisy reading of
capability (noise $\tau_K$), as one might obtain from proposal evaluation and expert
judgment. The review signal's informativeness is governed by $\tau_K$: the smaller the
noise, the more informative the review. Bayesian strategies update beliefs about every
researcher after each round and allocate on posterior expected marginal returns.

We compare seven funding strategies, which differ along two dimensions: information, with
or without the review signal, and planning, myopic or forward-looking. A myopic strategy
optimizes expected output in the current round; a forward-looking strategy optimizes
expected output over the whole remaining horizon. Three non-Bayesian baselines complete
the set: no funding, against which every result is normalized; uniform funding, which
splits each round's tranche equally; and track-record funding, which allocates in
proportion to the previous round's output. Uniform seed grants and lotteries, which modify
these strategies rather than add to them, are introduced in \S8.

\begin{table}[t]
\centering
\footnotesize
\setlength{\tabcolsep}{4pt}
\begin{tabular*}{\textwidth}{@{\extracolsep{\fill}}lll@{}}
\toprule
Strategy & Information & Planning \\
\midrule
No funding & none & none \\
Uniform & none & none \\
Track record & track record & none \\
Myopic, without review & track record + resource signal & current round \\
Forward, without review & track record + resource signal & remaining horizon \\
Myopic, with review & track record + resource signal + review & current round \\
Forward, with review & track record + resource signal + review & remaining horizon \\
\bottomrule
\end{tabular*}
\caption{The funding strategies compared.}
\label{tab:strategies}
\end{table}

Our conclusions rest either on analytic results supported by proofs or on simulation
results that are established across the following parameter ranges.\footnote{Unless a figure
states otherwise, unstated parameters take the values $T = 2$, $n = 50$, $\epsilon = 0.1$,
$b = 0.5$, $\alpha_K = \alpha_R = 2$, $\tau_K = \tau_R = 1$, and $\rho = 0$.
% TODO: move this note to the first figure's caption once figures land.
} The tail parameters $\alpha_K$ and $\alpha_R$ range from 1.3 to 3.5. Distributions of
research productivity are heavy-tailed: in Lotka's classic data, the number of authors
producing $x$ papers falls off as $x^{-2}$, and later studies estimate exponents mostly
between 2 and 3 \citep{Lotka1926, Nicholls1989}, % [check Nicholls range]
corresponding to Pareto tail parameters near 1 to 2; our range covers this territory and
extends to markedly more equal fields. The budget scale $b$ ranges from 0.1 to 1, from a
funder whose purse is small beside the community's own resources to one able to relieve
most researchers' bottlenecks; \S5 and \S7 also examine deeper purses. The compounding
rate $\epsilon$ ranges from near zero, where funding raises only current output, to 0.85,
where this round's funding substantially raises later capability. The review noise
$\tau_K$ ranges from 0.05, nearly perfect review, to 20, nearly uninformative review, a
range wide enough to contain the noise implied by empirical estimates of review's
predictive accuracy (\S6). The correlation $\rho$ ranges from $-0.5$ to 0.8, and the
horizon $T$ from one round to ten.

Randomness enters our model through the initial population draws, the signals, and
realized output. Realized output shapes the funder's beliefs and, through them, its
allocations. Capability growth, by contrast, follows expected output: the update above
uses $\Lambda$ itself, so a researcher's accumulation depends on their capability and
resources, not on luck in any given round. This is a modeling choice: it treats a round as
containing enough separate acts of research for luck to average out of accumulation, while
leaving the funder's inference fully exposed to noise. For the same reason, we score each
strategy by the sum of expected outputs along its run, which equals realized output
averaged over repetitions of the same allocations; this removes simulation noise from
comparisons without favoring any strategy. Each configuration runs on 200 independently
drawn populations, with common random numbers across strategies, so every comparison
between strategies is paired. Throughout, we report each strategy's improvement over no
funding, in percent. The full specification, parameter table, and code are in Appendix
[X]. % TODO: same owed simulation-specification appendix as \S3's normalization pointer

With the model in place, the funder's problem is exact: choose grants, given beliefs, to
maximize expected total output over the horizon. Next we solve the case in which the
funder has complete information regarding researchers' capabilities and resources.


\section{The optimal allocation: fund the capability-resource gap}\label{sec:optimal}

First, consider the simple case where the funder has complete information regarding
researchers' capabilities and resources. Further, we initially restrict our attention to a
single round of funding, temporarily putting aside the contribution of resources to the
accumulation of researchers' capabilities across rounds. The funder's problem is to divide
the round's tranche among researchers so as to maximize expected output:
\[
\max_{g_1, \ldots, g_n \geq 0} \; \sum_{i} A\,\Lambda(K_i,\, R_{i0} + g_i)
\quad \text{subject to} \quad \sum_{i} g_i \leq \text{the round's tranche}.
\]

Now, consider the marginal value of a dollar for such a funder. Under the harmonic
production function,
\[
\frac{\partial \lambda_i}{\partial g_i} = \frac{2A\,K_i^2}{(K_i + R_{i0} + g_i)^2}.
\]
The marginal value of funding researcher $i$ increases with their capability and decreases
with the resources they already hold, including grants. An allocation is optimal exactly
when no dollar can be moved to a researcher who can produce more value with it, that is,
when marginal values are equal across all funded researchers and no unfunded researcher's
marginal value exceeds theirs.

\begin{prop}[the gap rule]\label{prop:gap-rule}
For every researcher, the optimal grant fills the gap, if any, between the resources they
hold and a target proportional to their capability:
\[
g_i^* = \max(c\,K_i - R_{i0},\, 0),
\]
where the constant $c > 0$ is set by spending the budget.\footnote{$c = \sqrt{2A/\nu} - 1$,
with $\nu$ the common marginal value at the optimum. The proof is in Appendix~\ref{app:gap-rule}.}
\end{prop}

We can understand the rule as follows. Each researcher's target is $c$ times their
capability, and their gap is the difference between their target and the resources they
hold. The rule grants every researcher with a positive gap exactly their gap, and grants
nothing to the rest; the constant $c$ is defined as the value at which these grants sum to
the budget. Because $c$ is defined by the budget, the budget always covers the gaps: a
question of which gaps to fund first does not arise. When the budget is small, $c$ is
small: targets are low, few researchers have positive gaps, and their grants are small.
When the budget is large, $c$ is large: targets are high, and funding extends across
nearly the whole population. In both cases the funded researchers, those with positive
gaps ($R_{i0} < c\,K_i$), are those with the least resources relative to their capability,
and these are the researchers for whom a dollar has the greatest marginal
value.\footnote{The optimality condition behind the rule, funding to a common marginal
value, is the water-filling solution familiar from information theory
\citep{CoverThomas2006}. If all researchers shared one capability, the rule would reduce
to filling each researcher's gap below one common target, the poverty-gap transfer of
development economics \citep{FGT1984}; our rule scales the target with capability.} The
condition $R_{i0} < c\,K_i$ draws a frontier through the population, the funding frontier,
separating the funded from the unfunded (Figure~\ref{fig:frontier}).

% Figure 1: the gap rule in one population (R, K space, funding frontier, grants as arrows).
% Same population as Figure 2 (n = 75, K, R0 ~ Pareto(alpha = 2, min 1), rho = 0, A = 1,
% numpy default_rng(659)); budget scale b = 0.1, giving c = 0.833. Data: fig1-funded.dat
% (funded researchers with grant vectors), fig1-unfunded.dat. Generated by the snippet in
% state.md (2026-09-05); grants computed from the budget identity by bisection.
\begin{figure}[t]
\centering
\begin{tikzpicture}
\begin{axis}[
  width=0.82\textwidth, height=7.2cm,
  axis lines*=left,
  xmin=0, xmax=8.6, ymin=0, ymax=6.5,
  xtick=\empty, ytick=\empty,
  x axis line style={-{Stealth[length=4pt]}, line width=0.4pt},
  y axis line style={-{Stealth[length=4pt]}, line width=0.4pt},
  xlabel={resources $R$},
  ylabel={capability $K$},
  ylabel style={font=\small}, xlabel style={font=\small},
  tick label style={font=\small},
  axis line style={line width=0.4pt},
  tick style={line width=0.4pt, black},
  clip=false,
]
% funding frontier K = R/c
\addplot[black, line width=0.5pt, domain=0:6.4, samples=2] ({x*0.8332}, {x});
\node[font=\small, anchor=west, align=left] at (axis cs:5.6,6.25)
  {funding frontier ($R = cK$)};
% grants: horizontal arrows from (R0, K) to (cK, K)
\addplot[quiver={u=\thisrow{u}, v=\thisrow{v}}, -{Stealth[length=3.5pt]},
         gray, line width=0.5pt] table {fig1-funded.dat};
% funded researchers (filled)
\addplot[only marks, mark=*, mark size=1.6pt, black] table[x=R, y=K] {fig1-funded.dat};
% unfunded researchers (open)
\addplot[only marks, mark=o, mark size=1.6pt, gray, line width=0.5pt]
  table[x=R, y=K] {fig1-unfunded.dat};
% direct labels
\node[font=\small, anchor=west] at (axis cs:0.4,5.2) {funded};
\node[font=\small, anchor=west] at (axis cs:6.0,2.6) {unfunded};
\end{axis}
\end{tikzpicture}
\caption{The gap rule in one population. Each point is a researcher, placed by resources and capability; the line is the funding frontier $R = cK$. Researchers below the frontier (filled) are funded, and each arrow is a grant, which moves its researcher exactly to the frontier. Researchers on or beyond the frontier (open) receive nothing.}
\label{fig:frontier}
\end{figure}


The gap rule explains why each of the intuitive funding schemes discussed earlier fails.
Funding the track record fails because output reflects resources as well as capability: a
productive researcher whose resources approach or exceed their target has a small or
negative gap, and an additional dollar adds little output there. The failure is not merely a
shortfall from the optimum: because track-record funding channels money toward researchers
whose bottleneck is not resources, it can produce less output than spreading the budget
evenly, that is, than not targeting at all.\footnote{Example~\ref{ex:track-record} in Appendix~\ref{app:gap-rule} makes this exact with two researchers.} Funding the under-resourced fails for the mirror-image reason: scarcity is no
evidence of capability, and dollars given to researchers of low capability add little
output however scarce their resources. It, too, can produce less than not targeting at
all.\footnote{Again, Example~\ref{ex:scarcity} in Appendix~\ref{app:gap-rule} makes this exact.} Each scheme accounts for one of
capability and resources, whereas the marginal value of a dollar is determined by how a
researcher's resources compare with their capability, and so responds to both.

How much targeting matters depends on the budget. A researcher's expected output increases
with every dollar of resources, but with diminishing marginal returns, and it never
exceeds $2AK_i$: expected output is bounded above by a quantity set by capability, which
it approaches but does not attain. When the budget is large, optimal and uniform funding
therefore differ little in the output they produce: under either, every researcher's
output is near its upper bound, and the difference between the two vanishes as the budget
increases (Corollary~\ref{cor:targeting-vanishes}, Appendix~\ref{app:gap-rule}). When the budget is small, dollars are scarce, researchers differ
in the marginal value of a dollar, and output depends on where each dollar goes. Targeting
matters where money is tight and little where it is ample
(Figure~\ref{fig:targeting-value}). This comparative static recurs: it sets how much output is
lost by overriding the optimal allocation with seed grants and lotteries (\S8), and it determines when
optimal funding concentrates rather than spreads (\S8).

% Figure 2: the value of targeting decreases in the budget.
% Computed exactly from Proposition 1 (no simulation): population of n = 75 with
% K, R0 ~ Pareto(alpha = 2, min 1), rho = 0, A = 1, seed 659 (numpy default_rng);
% value of targeting = (f_opt - f_unif) / (f_unif - f_0), in percent;
% b = budget / total baseline resources. Data: fig2-curve.dat, generated by
% verify_gap_rule.py's companion snippet (see state.md, 2026-09-05).
\begin{figure}[t]
\centering
\begin{tikzpicture}
\begin{axis}[
  width=0.82\textwidth, height=5.4cm,
  axis lines*=left,
  xmin=0, xmax=2, ymin=0, ymax=175,
  xtick={0, 0.5, 1, 1.5, 2},
  ytick={0, 50, 100, 150},
  xlabel={budget scale $b$},
  ylabel={value of targeting (\%)},
  ylabel style={font=\small}, xlabel style={font=\small},
  tick label style={font=\small},
  axis line style={line width=0.4pt},
  tick style={line width=0.4pt, black},
  clip=false,
]
\addplot[black, line width=0.9pt, smooth] table {fig2-curve.dat};
\end{axis}
\end{tikzpicture}
\caption{The value of targeting decreases in the budget. The value of targeting is
the additional expected output from the optimal allocation over uniform funding, as
a fraction of uniform funding's gain over no funding; the budget scale $b$ is the
budget as a fraction of the population's total baseline resources. The curve shows
one population of $n = 75$ researchers drawn from the model's default distributions
($\alpha_K = \alpha_R = 2$, $\rho = 0$, $A = 1$); the value of targeting decreases
toward zero as the budget increases (Corollary~\ref{cor:targeting-vanishes}).}
\label{fig:targeting-value}
\end{figure}


To use the gap rule, however, a funder requires an estimate of the gap between
researchers' capabilities and their resources. Next we turn to how such estimation can
work: first from the track record alone, then with peer review added.

\section{What the track record can and cannot supply}\label{sec:records}

The funder observes track records: each researcher's realized output, round by round. A
track record is evidence about the capability-resource gap, and the question for this
section is how far that evidence goes.

Consider first what output can identify in principle. Expected output is one number,
produced by capability and resources jointly: every pair of capability and resources with
$2AKR/(K+R) = \lambda$ produces the same expected output $\lambda$. A researcher of high
capability and modest resources and a researcher of modest capability and abundant
resources can therefore produce identical records. Output does carry some information
about capability, a lower bound: since expected output never exceeds $2AK$, a researcher
who produces $\lambda$ has capability above $\lambda/2A$. Beyond that bound, the record
cannot separate the two researchers above, and these are exactly the researchers the gap
rule treats most differently: the first may have the largest gap in the population, the
second may have none. Cross-sectional output, however carefully recorded, underdetermines
what the funder most needs to know.

Our model measures how much output this underdetermination loses. A Bayesian funder relying on
records alone allocates far from the gap rule, and further rounds of records improve its
allocation only slowly; its output falls well short of the complete-information
benchmark.\footnote{At the default parameters, the correlation between its grants and
the optimal grants is 0.08 in the first round and reaches only 0.41 by round twenty; its
output falls short of the complete-information benchmark by 11.4 percent of no-funding
output (24.6 percent under heavy-tailed capability).} This shortfall is the span from
records alone to complete
information. Next we ask how much of that span a review signal recovers, and there
display this records-only baseline beside a funder holding the review signal.

The record is least informative exactly where the gap is largest. When a researcher's
resources are small relative to their capability, expected output is approximately
$2AR$: proportional to resources, nearly independent of capability
(Figure~\ref{fig:thin-grants}). On a thin grant, a researcher of capability 1 and a
researcher of capability 100 produce nearly indistinguishable output. Small grants
therefore yield almost no information about who is capable; separating capability from
resources by observing output requires funding depth, grants comparable to capability
itself. In our simulations of a resource-poor community with no review signal, a funder
allocating at standard depth gains almost nothing over uniform funding, while the same
funder with deep grants recovers most of the value of discrimination (\S7 measures this
and its consequences for timing). For nascent and resource-poor fields the implication
is direct: thin grants cannot reveal who is capable.

% Figure (proposal, for S5): small grants are uninformative about capability.
% Analytic curves lambda = AKR/(K+R) (drawn with A = 2, schematic axes), K in {3, 10, 30}; tangent line AR.
% Every curve has slope A at R = 0, so all coincide near the origin. Schematic axes
% per the Fig 1 convention (exact values do not matter). Data: fig3-curves.dat.
\begin{figure}[t]
\centering
\begin{tikzpicture}
\begin{axis}[
  width=0.82\textwidth, height=6.0cm,
  axis lines*=left,
  xmin=0, xmax=6.6, ymin=0, ymax=10.8,
  xtick=\empty, ytick=\empty,
  x axis line style={-{Stealth[length=4pt]}, line width=0.4pt},
  y axis line style={-{Stealth[length=4pt]}, line width=0.4pt},
  xlabel={resources $R$},
  ylabel={expected output},
  ylabel style={font=\small}, xlabel style={font=\small},
  axis line style={line width=0.4pt},
  clip=false,
]
\addplot[black, line width=0.7pt] table[x=R, y=k30] {fig3-curves.dat};
\addplot[black, line width=0.7pt] table[x=R, y=k10] {fig3-curves.dat};
\addplot[black, line width=0.7pt] table[x=R, y=k3]  {fig3-curves.dat};
\node[font=\small, anchor=west] at (axis cs:6.7,10.0) {$K = 30$};
\node[font=\small, anchor=west] at (axis cs:6.7,7.5)  {$K = 10$};
\node[font=\small, anchor=west] at (axis cs:6.7,4.0)  {$K = 3$};
\end{axis}
\end{tikzpicture}
\caption{On small grants, research output carries little information about capability. Expected output against resources for researchers of capability $K = 3$, $10$, and $30$. Where resources are small relative to capability the curves coincide, so researchers of very different capability produce nearly identical output; the curves separate only where resources are comparable to capability.}
\label{fig:thin-grants}
\end{figure}


What the record cannot supply, then, is a signal of capability separate from resources.
That is what peer review can provide under the right conditions. Next we measure what
such a signal is worth, and state the conditions.

\section{Peer review: what a capability signal is worth}\label{sec:review}

Peer review enters our model as a second source of evidence about researchers: a noisy
signal of each researcher's capability. The value of review, for a Bayesian funder, is
the expected value of its reduction of the funder's uncertainty regarding the
capability-resource gap. We can quantify this value in terms of output: the additional
expected output of a funder allocating with the review signal over the same funder
allocating without it.

The direction of review's effect is what one would expect: a more informative signal
brings the funder's allocation closer to the optimal allocation and its output closer
to what complete information would deliver, and with a sharp enough signal the funder
recovers most of the output that a funder relying on records alone leaves
unrealized.\footnote{At the default parameters, the correlation between the funder's
grants and the optimal grants increases from 0.22 at the noisiest signal in our range
to 0.94 at the sharpest; under heavy-tailed capability ($\alpha_K = 1.3$), from 0.34 to
0.96. The output shortfall to the complete-information benchmark decreases from 10.4 to
1.6 percent of no-funding output; under heavy-tailed capability, from 17.2 to 1.7
percent. The increase in review's value is not strictly monotone at the sharpest
signals: under heavy-tailed capability, its value at $\tau_K = 0.05$ is about two
percent below its value at $\tau_K = 0.3$, a small but statistically solid reversal
(paired across seeds, $z = 4.5$). The cause is the funder's other estimate: resources
are also observed with noise, and a slightly noisy capability signal tempers the
allocation against errors in the resource estimate, while a near-perfect one commits
the allocation to them. With the resource signal nearly noiseless, the sharper
capability signal is again worth more (paired $z = 5.8$).} Two further effects are less
obvious: (1) a funder relying on records alone does not catch up, and (2) what a review
signal is worth varies widely with the field. The first effect has a structural cause.
Capabilities compound across rounds while resources do not accumulate, so each
researcher's output approaches the level their resources sustain, and output at that
level carries less and less information about capability. Further rounds of records do
improve the funder's allocation, but slowly: after twenty rounds of records the
funder's allocation remains further from the optimal allocation than a funder with the
review signal is after one (Figure~\ref{fig:records-rounds}).

% Figure 5 (S5): records improve the allocation only slowly across rounds.
% y = shortfall of the funder's expected output below the complete-information optimum,
%     as a share of that optimum's gain over no funding, by round; T = 20; default
%     parameters, tau_K = 1, 50 seeds. No-review = S4; with review = S5.
% Data: fig5-rounds.dat (sf4, sf5; the s4, s5 columns hold the grant correlations
%     quoted in the text: 0.08 -> 0.32 vs 0.83), from byround_m.R at M = 1000 (2026-09-29).
% Heavy-tailed counterparts (alpha_K = 1.3, mean 2): shortfall 0.61 -> 0.16 vs 0.07 -> 0.03.
\begin{figure}[t]
\centering
\begin{tikzpicture}
\begin{axis}[
  width=0.82\textwidth, height=6.0cm,
  axis lines*=left,
  xmin=1, xmax=20, ymin=0, ymax=0.5,
  xtick={1, 5, 10, 15, 20},
  ytick={0, 0.1, 0.2, 0.3, 0.4, 0.5},
  yticklabels={0, 0.1, 0.2, 0.3, 0.4, 0.5},
  xlabel={round},
  ylabel={shortfall (share of complete-information gain)},
  ylabel style={font=\small}, xlabel style={font=\small},
  tick label style={font=\small},
  axis line style={line width=0.4pt},
  tick style={line width=0.4pt, black},
  clip=false,
]
\addplot[black, line width=0.7pt] table[x=round, y=sf4] {fig5-rounds.dat};
\addplot[black, line width=0.7pt] table[x=round, y=sf5] {fig5-rounds.dat};
\node[font=\small, anchor=west] at (axis cs:20.4,0.04) {with review};
\node[font=\small, anchor=west] at (axis cs:20.4,0.20) {no review};
\end{axis}
\end{tikzpicture}
\caption{The shortfall below complete information, round by round. The shortfall of a funder's expected research output below the complete-information optimum, as a share of what complete-information funding adds, over twenty rounds in the intermediate field, for a funder without review and a funder with one review score observed at the start. Without review the shortfall is 46 percent in round one and 19 percent in round twenty; with review, 8 percent and 3 percent. Settings: Appendix~\ref{app:simulation}.}
\label{fig:records-rounds}
\end{figure}


Our main result delineates how much review is worth as a function of the field
(Figure~\ref{fig:review-value}). Review is worth most where capability is heavy-tailed
and the budget is tight, and in the heavy-tailed case even a signal of moderate
informativeness captures most of its value. The reason is twofold. Where capability is
heavy-tailed, most of the output a funder can add comes from getting resources to the
few researchers of unusually high capability. And precisely those researchers are the
easiest to identify: they stand far apart from everyone else, so even a noisy signal
separates them. Identifying the most capable researchers captures most of the value and
requires the least information. Sharpening review beyond a modest level therefore adds
little.\footnote{At three times the default signal noise, review retains 82 percent of
its maximum value under heavy-tailed capability ($\alpha_K = 1.3$), 44 percent at the
default distribution ($\alpha_K = 2$), and 10 percent where capability is spread more
evenly ($\alpha_K = 3.5$). The noise level at which review's value falls to half its
maximum is roughly $\tau_K = 10$, $2.5$, and $0.75$ in the three cases; the default
noise is $\tau_K = 1$, and our sweep runs from near-perfect review ($\tau_K = 0.05$) to
nearly uninformative review ($\tau_K = 20$). Fifty seeds per cell.} This ordering of
review's value by field does not depend on our production form: it holds across the
family of technologies from Cobb-Douglas to Leontief (Appendix~\ref{app:technology}).

% Figure 4 (S5): the value of review by field.
% y = share of the no-review shortfall that review recovers, (S5 - S4)/(oracle - S4),
%     mean and standard error over 100 simulated populations per point;
% x = review noise tau_K, log scale, 0.05 to 20 (default 1).
% Capability Pareto with tail alpha_K in {1.3, 2, 3.5} and the SCALE SET SO THAT E[K] = 2
% FOR EVERY TAIL (k_min = 2(alpha-1)/alpha), so the sweep varies inequality alone
% (revision 2026-09-29, referee point 2). Data: fig4-value.dat, from sens2.R (norm_fine),
% model.R with the smooth allocator, T = 2, b = 1.
\begin{figure}[t]
\centering
\begin{tikzpicture}
\begin{axis}[
  width=0.82\textwidth, height=6.0cm,
  axis lines*=left,
  xmode=log,
  xmin=0.05, xmax=20, ymin=0, ymax=1,
  xtick={0.05, 0.3, 1, 3, 10, 20},
  xticklabels={0.05, 0.3, 1, 3, 10, 20},
  ytick={0, 0.25, 0.5, 0.75, 1},
  yticklabels={0, 0.25, 0.5, 0.75, 1},
  xlabel={signal noise $\tau_K$},
  ylabel={share of shortfall recovered},
  ylabel style={font=\small}, xlabel style={font=\small},
  tick label style={font=\small},
  axis line style={line width=0.4pt},
  tick style={line width=0.4pt, black},
  clip=false,
]
\addplot[black, line width=0.7pt, error bars/.cd, y dir=both, y explicit, error bar style={line width=0.4pt}, error mark options={rotate=90, mark size=1.2pt, line width=0.4pt}] table[x=tau, y=s13, y error=e13] {fig4-value.dat};
\addplot[black, line width=0.7pt, error bars/.cd, y dir=both, y explicit, error bar style={line width=0.4pt}, error mark options={rotate=90, mark size=1.2pt, line width=0.4pt}] table[x=tau, y=s20, y error=e20] {fig4-value.dat};
\addplot[black, line width=0.7pt, error bars/.cd, y dir=both, y explicit, error bar style={line width=0.4pt}, error mark options={rotate=90, mark size=1.2pt, line width=0.4pt}] table[x=tau, y=s35, y error=e35] {fig4-value.dat};
\node[font=\small, anchor=west] at (axis cs:22,0.13)  {$\alpha_K = 1.3$};
\node[font=\small, anchor=west] at (axis cs:22,0.035) {$\alpha_K = 2$};
\node[font=\small, anchor=west] at (axis cs:22,-0.06) {$\alpha_K = 3.5$};
\end{axis}
\end{tikzpicture}
\caption{The value of review by field. The share of the no-review shortfall that review recovers, as review noise $\tau_K$ grows, in three fields with the same mean capability: heavy-tailed ($\alpha_K = 1.3$), intermediate ($2$), and evenly spread ($3.5$). Bars: one standard error over 100 populations. Settings: Appendix~\ref{app:simulation}.}
\label{fig:review-value}
\end{figure}


A corollary concerns where review effort goes. In our model, what adds output is separating
the exceptional researchers from the rest, and that separation requires little
information; fine distinctions among the many applications of comparable middling merit
add comparatively little output. A review system that invests most of its effort in
exactly those fine distinctions may be misallocating it.

The remaining comparisons run the other way. Where capability is spread more evenly, no
researcher matters much more than another, so review of any informativeness adds
little.\footnote{At $\alpha_K = 3.5$, the entire records-only shortfall is 3.5 percent
of no-funding output (24.6 percent under heavy-tailed capability), and even
near-perfect review recovers only 72 percent of it: review's value at its maximum is
2.5 percent of no-funding output, against 23 percent under heavy-tailed capability.}
And where the budget is ample, optimal and uniform funding produce nearly the same
output, so information about whom to fund has little left to add: review is worth most
where money is tight.

These results bear on the dispute over the efficacy of peer review. Among funded NIH
grants, percentile scores predict subsequent productivity barely better than chance. In
our model, a signal of that accuracy is very noisy review.\footnote{An AUC of 0.54
\citep{Fang2016} corresponds to $\tau_K \gtrsim 20$ at the default parameters, the
noisiest signal in our sweep. The mapping from measured accuracy to signal noise itself
depends on the field: at the same noise, accuracy is higher under heavy-tailed
capability, because the most capable researchers stand further apart.} Whether such
review is worth having depends on the field. At the noisiest signal in our sweep,
review still recovers about a third of the records-only shortfall under heavy-tailed
capability, under a tenth at the default distribution, and essentially none where
capability is spread more evenly.\footnote{31, 9, and 0.1 percent of the records-only
shortfall at $\tau_K = 20$, for $\alpha_K = 1.3$, $2$, and $3.5$; fifty seeds per
cell.} Accuracy is one input to review's value; the field's capability distribution and
the funder's budget are the others. No single number settles the dispute: review of the
same informativeness is worth several times more in one field than in another, and the
same measured accuracy corresponds to different informativeness in different fields.

One condition qualifies all of this: the funder must weight review at its true
informativeness. A funder that overtrusts review, treating a noisy signal as sharp, can
do worse than one that ignores review altogether: with moderately spread capability,
overtrust turns review's value negative; with heavy-tailed capability it reduces the
value without reversing it (Figure~\ref{fig:overtrust}).\footnote{Each negative cell
in Figure~\ref{fig:overtrust} is within two standard errors of zero, but all are
negative. The tenfold comparison, across true noise levels: a funder that overtrusts
tenfold (true $\tau_K = 3$, belief $0.3$) loses more than review's full calibrated
value at the default capability distribution, while a funder that undertrusts tenfold
(true $\tau_K = 0.3$, belief $3$) retains over half of it.} Overtrusting review loses more output than
undertrusting it.

% Figure 6 (S6): the value of review under misweighted trust.
% y = review's value (S5 - S4) as percent of no-funding output; x = the funder's
% believed noise (log scale); true noise tau_K = 3 (dotted rule). Two capability
% distributions. Error bars: one standard error, 200 seeds per cell.
% Data: fig6-overtrust.dat, from sweep_results/D_misspecified_trust (hash 21b0d9a).
\begin{figure}[t]
\centering
\begin{tikzpicture}
\begin{axis}[
  width=0.82\textwidth, height=6.0cm,
  axis lines*=left,
  xmode=log,
  xmin=0.1, xmax=10, ymin=-2, ymax=20,
  xtick={0.1, 0.3, 1, 3, 10},
  xticklabels={0.1, 0.3, 1, 3, 10},
  ytick={0, 5, 10, 15, 20},
  xlabel={believed noise $\tau_K$},
  ylabel={value of review (\% of no-funding output)},
  ylabel style={font=\small}, xlabel style={font=\small},
  tick label style={font=\small},
  axis line style={line width=0.4pt},
  tick style={line width=0.4pt, black},
  clip=false,
]
\addplot[gray, dashed, line width=0.4pt, domain=0.1:10] {0};
\draw[gray, dotted, line width=0.5pt] (axis cs:3,-2) -- (axis cs:3,20);
\node[font=\small, gray, anchor=west] at (axis cs:3.15,8) {true noise};
\addplot[black, line width=0.7pt, error bars/.cd, y dir=both, y explicit]
  table[x=belief, y=p13, y error=s13] {fig6-overtrust.dat};
\addplot[black, line width=0.7pt, error bars/.cd, y dir=both, y explicit]
  table[x=belief, y=p20, y error=s20] {fig6-overtrust.dat};
\node[font=\small, anchor=west] at (axis cs:11,13.1) {$\alpha_K = 1.3$};
\node[font=\small, anchor=west] at (axis cs:11,1.2)  {$\alpha_K = 2$};
\end{axis}
\end{tikzpicture}
\caption{Overtrusting review loses more output than undertrusting it. The value of review as a
percent of no-funding output, as the funder's believed signal noise varies while the
true noise is $\tau_K = 3$ (dotted rule); bars show one standard error over two
hundred seeds. Value peaks where belief matches truth. To the right of the peak,
undertrust forgoes part of the value; to the left, overtrust treats a noisy signal
as sharp, and at the default capability distribution ($\alpha_K = 2$) it takes the
value below zero: worse than ignoring review. Under heavy-tailed capability
($\alpha_K = 1.3$) overtrust reduces the value without reversing it.}
\label{fig:overtrust}
\end{figure}


Review can inform whom to fund. Next we turn to when to fund: how funding should be
spread across rounds, and how much that choice matters.

\section{The timing of funding}\label{sec:timing}

Funding in our model is spread over rounds. The funder holds a fixed budget for the
whole horizon, and between rounds two things move: researchers produce output, which
the funder observes, and capabilities compound through the research that output
represents. Grants themselves are consumed within their round. Up to now the funder
has spent its budget in equal parts each round; this section asks whether it should
not: whether to spend early, evenly, or late, and how much the choice matters.

The strongest argument for spending early concerns resource-poor fields. Researchers
without resources produce nothing; where nothing is produced, nothing is observed and
nothing compounds; so money held for later rounds produces neither evidence nor growth in
the meantime. On this argument a funder facing a poor field should spend at once, to
start the engine. The argument fails on inspection. A funder that spends everything in
the first round allocates before observing any output, exactly as blind as one that
spends everything in the last; in the extreme case of a field with no resources at
all, the two schedules produce identical expected output. What the argument does
establish is a complementarity: early spending enables observation and growth, and
late spending exploits what early spending revealed. That favors putting some money
early. It does not favor putting most of it early.\footnote{In a planner-free check
over a grid of two-block schedules executed by a funder re-deciding each round, the
best schedule that puts an above-even share in the early rounds loses to the even
schedule, by 1.6 percent of the even schedule's output at negligible compounding and
by 3.5 percent at the default compounding rate (six rounds, resource-poor community,
two hundred seeds).}

Our simulations bear this out and say which way the asymmetry runs: observation comes
first, money after. Wherever capabilities compound at more than a token rate, the
optimal schedule holds a below-even share in the first round and deploys the bulk once
output has revealed who is capable; spending late is rewarded because late money is
informed money, and because early output compounds on its own
(Figure~\ref{fig:timing-boundary}). Resource poverty raises the early share and
decreases the intensity of the effect, but does not reverse
it.\footnote{Figure~\ref{fig:timing-boundary} shows three of the five baseline
resource scales in the map; the other two lie between the curves shown. Every cell
whose schedule leans early has a compounding rate $\epsilon \leq 0.03$, and there the
schedule's center of mass is within 0.005 of even.}

% Figure (S7): the schedule's tilt is governed by the compounding rate.
% y = budget-schedule center of mass (0.5 = even; above = late-leaning) of the
% forward funder; x = compounding rate eps (log scale); three baseline resource
% scales (r_min = 0.001 poor, 0.3, 3 rich). 200 seeds per cell.
% Data: fig8-boundary.dat, from sweep_results/resource_regime (hash 21b0d9a).
\begin{figure}[t]
\centering
\begin{tikzpicture}
\begin{axis}[
  width=0.82\textwidth, height=5.4cm,
  axis lines*=left,
  xmode=log,
  xmin=0.0001, xmax=0.85, ymin=0.48, ymax=0.68,
  xtick={0.0001, 0.01, 0.1, 0.85},
  xticklabels={0.0001, 0.01, 0.1, 0.85},
  ytick={0.5, 0.55, 0.6, 0.65},
  yticklabels={even, 0.55, 0.60, 0.65},
  xlabel={compounding rate $\epsilon$},
  ylabel={schedule center of mass},
  ylabel style={font=\small}, xlabel style={font=\small},
  tick label style={font=\small},
  axis line style={line width=0.4pt},
  tick style={line width=0.4pt, black},
  clip=false,
]
\addplot[gray, dashed, line width=0.4pt, domain=0.0001:0.85] {0.5};
\addplot[black, line width=0.7pt] table[x=eps, y=rich] {fig8-boundary.dat};
\addplot[black, line width=0.7pt] table[x=eps, y=mid]  {fig8-boundary.dat};
\addplot[black, line width=0.7pt] table[x=eps, y=poor] {fig8-boundary.dat};
\node[font=\small, anchor=west] at (axis cs:0.92,0.652) {resource-rich};
\node[font=\small, anchor=west] at (axis cs:0.92,0.566) {intermediate};
\node[font=\small, anchor=west] at (axis cs:0.92,0.521) {resource-poor};
\end{axis}
\end{tikzpicture}
\caption{The compounding rate, not resource poverty, sets how late the schedule leans. The schedule's center of mass (0.5 is even; higher is later) as the compounding rate $\epsilon$ increases, for three baseline resource scales. Settings: Appendix~\ref{app:simulation}.}
\label{fig:timing-boundary}
\end{figure}


How much a funder can learn from its own grants depends on their size. A researcher on
a thin grant produces output nearly proportional to the grant and nearly independent
of capability, so a funder whose grants are thin learns almost nothing it can exploit
later: in a resource-poor field with no review signal, such a funder gains almost
nothing over uniform funding. Deep grants change this. As grants approach the scale of
capability itself, output separates the capable from the rest, and the funder recovers
most of the value of discrimination, on a schedule that puts observation before money:
a small early share to make output observable, the mass deployed once informed
(Figure~\ref{fig:depth-schedule}).\footnote{In a resource-poor community with
heavy-tailed capability, no review signal, and six rounds (two hundred seeds): the
informed funder's gain over uniform funding increases from a third of one percent of
uniform funding's output at standard depth to nine percent at six times that depth and
ten percent at twelve times, roughly the gain a sharp review signal delivers at
standard depth (nine to eleven percent). In these runs the budget is scaled against
capability rather than resources: at the six-times depth, the average grant per
researcher per round is half of mean capability.}

% Figure (S7): the optimal schedule at standard and deep funding.
% y = share of the budget spent in each round by the forward funder (S8); T = 6;
% poverty corner (r_min = 0.001), heavy tail (alpha_K = 1.3), no review signal
% (tau_K = 100), negligible compounding (eps = 1e-4), budget scaled against
% capability (budget_ref = "K"). b = 0.5: 50 seeds; b = 3: 24 seeds (max SE 0.006);
% round-1 share at b = 3 matches the canonical 200-seed run (0.104 vs 0.110).
% Data: fig7-schedule.dat, generated by for-claude/verify_s7.R (schedule addendum).
\begin{figure}[t]
\centering
\begin{tikzpicture}
\begin{axis}[
  width=0.82\textwidth, height=5.4cm,
  axis lines*=left,
  xmin=1, xmax=6, ymin=0.08, ymax=0.20,
  xtick={1, 2, 3, 4, 5, 6},
  ytick={0.10, 0.15, 0.167, 0.20},
  yticklabels={0.10, 0.15, even, 0.20},
  xlabel={round},
  ylabel={share of budget},
  ylabel style={font=\small}, xlabel style={font=\small},
  tick label style={font=\small},
  axis line style={line width=0.4pt},
  tick style={line width=0.4pt, black},
  clip=false,
]
\addplot[black, line width=0.7pt, mark=*, mark size=1.2pt] table[x=round, y=b05] {fig7-schedule.dat};
\addplot[black, line width=0.7pt, mark=o, mark size=1.4pt] table[x=round, y=b3] {fig7-schedule.dat};
\node[font=\small, anchor=west] at (axis cs:6.15,0.167) {standard depth ($b = 0.5$)};
\node[font=\small, anchor=west] at (axis cs:6.15,0.183) {deep grants ($b = 3$)};
\end{axis}
\end{tikzpicture}
\caption{Observation first, money after. The share of the budget the funder spends in
each of six rounds, in a resource-poor field with heavy-tailed capability and no
review signal. At standard depth (filled marks) the funder learns nothing from its
own thin grants and the optimal schedule is even. With deep grants (open marks) the
funder holds the first-round share below even and deploys the bulk once output has
revealed who is capable.}
\label{fig:depth-schedule}
\end{figure}


The pattern has a single main cause. Capabilities compound through all research
output, funded or not; grants add to output and so add to compounding. Separating the
two channels: with only the grant-fed channel active, the optimal schedule leans
mildly early, since earlier grants compound longer. With only the free channel active,
it leans strongly late. The free channel dominates once it reaches even a fraction of
the grant-fed rate, so in any field where research compounds regardless of who is
funded, the funder's reason to wait outweighs its reason to hurry.\footnote{With the
free channel off, the schedule's center of mass is 0.48 to 0.49 across grant-fed
rates, strictly below even; with the grant-fed channel off, 0.53 to 0.68. The free
channel's rate takes over at roughly one eighth to one third of the grant-fed rate
(two hundred seeds, five rounds).}

Finally, what is the schedule choice worth? Less than information, in every setting we
test. Across all horizons and compounding rates in our sweeps, the gain from
deliberate scheduling over even installments re-decided each round stays below the
review signal's value: at the default compounding rate and below it is smaller by an
order of magnitude or more, and it approaches the signal's value only at the joint
extreme of the highest compounding rate and the longest horizon we consider
(Figure~\ref{fig:schedule-vs-signal}).\footnote{As a fraction of no-funding output,
the scheduling gain ranges from indistinguishable from zero to 0.7 percent across the
thirty-two tested cells; the review signal's value in the same cells ranges from 0.8
to 7.5 percent.} The gain increases roughly with the square of the horizon over short
horizons and saturates beyond them.\footnote{Fitted exponents 2.3 and 2.1 for horizons
up to five rounds at $\epsilon = 0.3$ and $0.85$; 1.4 and 0.8 over five to ten
rounds.} And it exists only where capability and resources are complements, as in our
production form: under Cobb-Douglas production the value of scheduling is
indistinguishable from zero and the optimal schedule stays within about a hundredth of
even (Appendix~\ref{app:technology}, Table~\ref{tab:timing-technology}). For a funder thinking about how to maximize their impact, choosing whom to fund
outweighs choosing when.

% Figure (S7): what the schedule choice is worth, relative to the review signal.
% y = gain from deliberate scheduling (S8 - S5) as a share of the review signal's
% value (S8 - S7), by horizon T and compounding rate eps; negative gains floored at
% zero (all within noise of zero). Default parameters otherwise; 64-200 seeds per
% cell (T_run_smooth horizon_growth + horizon_long, hash 21b0d9a; re-derived
% in-session, verify_s7.R). Curves for eps in {0.05, 0.15, 0.3, 0.55, 0.85};
% eps = 0.05, 0.15, 0.55 were swept to T = 5 only.
\begin{figure}[t]
\centering
\begin{tikzpicture}
\begin{axis}[
  width=0.82\textwidth, height=5.7cm,
  axis lines*=left,
  xmin=2, xmax=10, ymin=0, ymax=1.05,
  xtick={2, 4, 6, 8, 10},
  ytick={0, 0.25, 0.5, 0.75, 1},
  yticklabels={0, 0.25, 0.5, 0.75, 1},
  xlabel={horizon $T$ (rounds)},
  ylabel={scheduling gain / signal value},
  ylabel style={font=\small}, xlabel style={font=\small},
  tick label style={font=\small},
  axis line style={line width=0.4pt},
  tick style={line width=0.4pt, black},
  clip=false,
]
\addplot[gray, dashed, line width=0.4pt, domain=2:10] {1};
\node[font=\small, gray, anchor=south west] at (axis cs:2.1,1.005) {equal to the signal's value};
\addplot[black, line width=0.7pt] coordinates {(2,0.0468) (3,0.1136) (4,0.1855) (5,0.2635) (6,0.3443) (7,0.4213) (8,0.4982) (9,0.5811) (10,0.6556)};
\addplot[black, line width=0.7pt] coordinates {(2,0.0198) (3,0.0482) (4,0.0769) (5,0.1090)};
\addplot[black, line width=0.7pt] coordinates {(2,0.0058) (3,0.0144) (4,0.0224) (5,0.0317) (6,0.0398) (7,0.0487) (8,0.0561) (9,0.0651) (10,0.0749)};
\addplot[black, line width=0.7pt] coordinates {(2,0.0012) (3,0.0034) (4,0.0051) (5,0.0069)};
\addplot[black, line width=0.7pt] coordinates {(2,0.0000) (3,0.0002) (4,0.0001) (5,0.0000)};
\node[font=\small, anchor=west] at (axis cs:10.2,0.656) {$\epsilon = 0.85$};
\node[font=\small, anchor=west] at (axis cs:5.15,0.126) {$\epsilon = 0.55$};
\node[font=\small, anchor=west] at (axis cs:10.2,0.075) {$\epsilon = 0.3$};
\end{axis}
\end{tikzpicture}
\caption{Choosing whom outweighs choosing when. The gain from deliberate scheduling
(a funder that plans its budget's spread over the horizon, over one that spends even
installments and re-decides each round) as a share of the review signal's value, by
horizon and compounding rate $\epsilon$; the default rate is $0.3$. In every setting
we test the gain is below the signal's value: an order of magnitude below or more at
the default rate and below, and approaching the signal's value only at the joint
extreme of the highest rate and the longest horizon. The two lowest curves, at $\epsilon = 0.15$ and $0.05$, and the curve at $0.55$ were
swept to five rounds.}
\label{fig:schedule-vs-signal}
\end{figure}


Next we turn to two alternatives to targeting that recur in funding policy: uniform
seed grants and lotteries.

\section{Spreading funds: seed grants, lotteries, and concentration}\label{sec:egalitarian}

Targeting sends each dollar to the researcher expected to produce the most with it.
Two alternatives to targeting recur in funding policy, and each has a rationale.
Uniform seed grants give every researcher the same amount, so that a researcher whom
review misjudges still receives something; in practice a funder gives out part of its
budget as seed grants and targets the rest. Lotteries give grants to applicants chosen
at random, usually from among those who pass a review threshold; their advocates
argue that review cannot reliably rank the applicants it deems fundable, that review
is costly, and that review is biased. Both alternatives give up part of targeting on
purpose. Our model can say how much expected output is lost by giving up targeting,
field by field, and where the loss is small. What the alternatives save, in the cost of
review and in the effort researchers spend on proposals, lies outside our model and
has been studied by others;\footnote{\citet{GrossBergstrom2019} model grant
competition as a contest in which scientists spend costly effort on proposals. They
show that when only a small fraction of proposals is funded, this effort can rival
the scientific value of the research funded, and that a lottery among proposals above
a quality threshold removes the link between the wasted effort and the fraction
funded.} we return to the comparison at the close of the section.

Both alternatives lose the same thing: the value of targeting, the additional expected
output that the optimal allocation produces over an equal division of the same
budget. How much output is lost by giving up targeting depends on the field. For a
funder allocating with records and review, targeting adds more than uniform funding's
entire gain over no funding in a heavy-tailed field with sharp review and a tight
budget, and about a quarter of uniform funding's gain in the default field at the
largest budget we sweep.\footnote{Over two rounds: the review-informed funder's gain
over uniform funding, as a fraction of uniform funding's gain over no funding,
decreases from 1.17 to 0.72 across budget scales $b = 0.1$ to $1$ with heavy-tailed
capability and sharp review ($\tau_K = 0.3$), and from 0.74 to 0.25 in the default
field with the default signal. The two regimes differ in both the capability
distribution and the signal's noise. Two hundred seeds per cell.} Giving up targeting
loses little output under the same conditions that make review worth little: evenly
spread capabilities, large budgets, uninformative peer review.

Consider uniform seed grants first. The output lost increases faster than the
fraction of the budget given out as seed grants: giving out half of the budget as
seed grants loses more than twice the output that giving out a quarter loses
(Figure~\ref{fig:floor-cost}). The size of the loss is set by what targeting is worth
in the field. When a review-informed funder gives out three quarters of its budget as
seed grants, expected output falls by about six percent of no-funding output in a
heavy-tailed field with sharp review, and by about one percent in the default
field.\footnote{Across the swept cells, the output lost is approximately the informed
funder's gain over uniform funding, times the fraction of the budget given out as
seed grants, times a factor between 0.10 and 0.86; the factor is above 0.3 at tight
budgets and in heavy-tailed fields. Two hundred seeds per cell.} When the whole
budget is given out as seed grants, the scheme is uniform funding, and the full value
of targeting is lost. Whether the reasons for seed grants justify the loss is a
judgment our model does not make; the model supplies the size of the loss.

% Figure (S8): the output lost to uniform seed grants.
% y = output lost by flooring the review-informed forward funder (S11 vs S8) as
%     percent of no-funding output; x = floored share of the budget (x_seed).
% Curves: heavy-tailed field with reliable review (tau_K = 0.3) at b = 1, 0.5, 0.1;
% default field with default review (tau_K = 1) at b = 0.5. T = 2, 200 seeds.
% Data: fig10-floor-cost.dat, from sweep_results/D3_seed_signal (hash 21b0d9a),
% re-derived in-session (verify_s8.R). Regime pairing caveat: the heavy-tailed cells
% run with sharper review than the default cells; stated in the caption.
\begin{figure}[t]
\centering
\begin{tikzpicture}
\begin{axis}[
  width=0.76\textwidth, height=5.4cm,
  axis lines*=left,
  xmin=0, xmax=0.75, ymin=0, ymax=20,
  xtick={0, 0.25, 0.5, 0.75},
  ytick={0, 5, 10, 15, 20},
  xlabel={fraction of the budget given out as seed grants},
  ylabel={research output lost (\% of the gain from funding)},
  ylabel style={font=\small}, xlabel style={font=\small},
  tick label style={font=\small},
  axis line style={line width=0.4pt},
  tick style={line width=0.4pt, black},
  clip=false,
]
\addplot[black, line width=0.7pt, mark=*, mark size=1.2pt] table[x=xseed, y=hb1]  {fig10-floor-cost.dat};
\addplot[black, line width=0.7pt, mark=*, mark size=1.2pt] table[x=xseed, y=hb05] {fig10-floor-cost.dat};
\addplot[black, line width=0.7pt, mark=*, mark size=1.2pt] table[x=xseed, y=hb01] {fig10-floor-cost.dat};
\addplot[black, line width=0.7pt, mark=*, mark size=1.2pt] table[x=xseed, y=db05] {fig10-floor-cost.dat};
\node[font=\small, anchor=west] at (axis cs:0.77,8.11) {heavy-tailed, $b = 2$};
\node[font=\small, anchor=west] at (axis cs:0.77,11.07) {heavy-tailed, $b = 1$};
\node[font=\small, anchor=west] at (axis cs:0.77,17.29) {heavy-tailed, $b = 0.2$};
\node[font=\small, anchor=west] at (axis cs:0.77,3.47) {default, $b = 1$};
\end{axis}
\end{tikzpicture}
\caption{The research output lost to uniform seed grants. Expected research output lost, as a percent of the research output that the funder's grants add over no funding, when a review-informed funder gives out a fraction of its budget as equal grants to all researchers and targets the rest. The loss grows faster than the fraction, and its size is set by what targeting is worth in the field. Settings: Appendix~\ref{app:simulation}.}
\label{fig:floor-cost}
\end{figure}


Now consider lotteries. A lottery gives each applicant in its pool a chance at a full
grant instead of a certain smaller one, and on average the applicant receives the
same amount either way. A researcher's expected output increases with resources at a
diminishing rate, so a certain grant produces more expected output than a chance at a
larger grant of the same average size. Dividing the pool's money equally among the
pool therefore produces more expected output than any lottery over the same pool, and
a lottery over all applicants produces less expected output than uniform funding
(Proposition~\ref{prop:lottery}, Appendix~\ref{app:lottery}). To measure how much
less, and how lotteries compare with selection by review, we evaluate five schemes in
a single round (Figure~\ref{fig:lottery-schemes}). The first scheme divides the budget
equally among the fifth of researchers with the highest review scores. The second
divides the budget equally among the top two fifths. The third is the partial lottery
that some funders run: it funds half of the top two fifths at random, with the same
grant size as the first scheme. The fourth is uniform funding. The fifth is a lottery
over all researchers, again with the first scheme's grant size. In a heavy-tailed
field with a tight budget, selection by review, the first scheme, gains about three
quarters of what the optimal allocation gains when review is sharp and about half
when review is very noisy. The partial lottery gains less than selection by review,
and most of the difference comes from ignoring the review ranking within the top two
fifths, not from the randomness itself.\footnote{At review noise $\tau = 1$: selection
by review gains 0.75 of what the optimal allocation gains, the equal division among
the top two fifths 0.61, and the partial lottery 0.56. Of the 0.18 by which the
partial lottery falls short of selection by review, ignoring the ranking accounts for
0.14 and the randomness for 0.04. Single round, a fifth of the population funded, two
thousand populations per point.} The lottery over all researchers gains least, less
than uniform funding. In an evenly spread field with a large budget, the ordering
reverses: uniform funding gains more than four fifths of what the optimal allocation
gains, and every scheme that concentrates grants, whether by review score or by
chance, gains less.

% Figure (S8): what spreading and randomizing capture, by field.
% y = a scheme's gain over no funding as a share of the optimal (complete-information
%     gap-rule) allocation's gain; x = screen noise tau (log). Single round, n = 50,
%     funded share one fifth of the population, 2000 populations per point (max MC SE
%     0.002); screen s = K + N(0, tau). Schemes: equal division among the top fifth
%     by screen, equal division among the top two fifths (half-size grants), partial lottery
%     (half of the top two fifths at random, full-size grants), uniform, lottery over all
%     (a fifth of the population at random). Left panel: heavy-tailed capability
%     (alpha_K = 1.3), tight budget (b = 0.1). Right panel: evenly spread capability
%     (alpha_K = 3.5), b = 0.5. Data: fig11-heavy.dat, fig11-even.dat
%     (price_lotteries.py, 2026-09-08).
\begin{figure}[t]
\centering
\begin{tikzpicture}
\begin{axis}[
  name=left,
  width=0.46\textwidth, height=5.6cm,
  axis lines*=left,
  xmode=log,
  xmin=0.3, xmax=10, ymin=0, ymax=1,
  xtick={0.3, 1, 3, 10},
  xticklabels={0.3, 1, 3, 10},
  ytick={0, 0.25, 0.5, 0.75, 1},
  yticklabels={0, 0.25, 0.5, 0.75, 1},
  xlabel={review noise $\tau$},
  ylabel={share of the optimal allocation's gain},
  ylabel style={font=\small}, xlabel style={font=\small},
  tick label style={font=\small},
  axis line style={line width=0.4pt},
  tick style={line width=0.4pt, black},
  title={\small heavy-tailed, tight budget},
  title style={yshift=-2pt},
  clip=false,
]
\addplot[gray, dashed, line width=0.5pt, domain=0.3:10] {0.445};
\addplot[gray, dotted, line width=0.6pt, domain=0.3:10] {0.378};
\addplot[black, line width=0.7pt, mark=*, mark size=1.1pt] table[x=tau, y=ss] {fig11-heavy.dat};
\addplot[black, line width=0.7pt, mark=o, mark size=1.3pt] table[x=tau, y=ws] {fig11-heavy.dat};
\addplot[black, line width=0.7pt, mark=triangle*, mark size=1.4pt] table[x=tau, y=sl] {fig11-heavy.dat};
\node[font=\scriptsize, anchor=south west] at (axis cs:0.315,0.79) {top fifth by review};
\node[font=\scriptsize, gray, anchor=north west] at (axis cs:0.315,0.372) {lottery over all};
\end{axis}
\begin{axis}[
  at={(left.outer east)}, anchor=outer west, xshift=6pt,
  width=0.46\textwidth, height=5.6cm,
  axis lines*=left,
  xmode=log,
  xmin=0.3, xmax=10, ymin=0, ymax=1,
  xtick={0.3, 1, 3, 10},
  xticklabels={0.3, 1, 3, 10},
  ytick={0, 0.25, 0.5, 0.75, 1},
  yticklabels={,,,,},
  xlabel={review noise $\tau$},
  xlabel style={font=\small},
  tick label style={font=\small},
  axis line style={line width=0.4pt},
  tick style={line width=0.4pt, black},
  title={\small evenly spread, ample budget},
  title style={yshift=-2pt},
  clip=false,
]
\addplot[gray, dashed, line width=0.5pt, domain=0.3:10] {0.831};
\addplot[gray, dotted, line width=0.6pt, domain=0.3:10] {0.414};
\addplot[black, line width=0.7pt, mark=*, mark size=1.1pt] table[x=tau, y=ss] {fig11-even.dat};
\addplot[black, line width=0.7pt, mark=o, mark size=1.3pt] table[x=tau, y=ws] {fig11-even.dat};
\addplot[black, line width=0.7pt, mark=triangle*, mark size=1.4pt] table[x=tau, y=sl] {fig11-even.dat};
\node[font=\scriptsize, gray, anchor=south west] at (axis cs:0.32,0.836) {uniform};
\node[font=\scriptsize, anchor=south west] at (axis cs:1.15,0.705) {top two fifths by review};
\end{axis}
\end{tikzpicture}
\caption{What five funding schemes gain, by field. Each scheme's gain over no funding
as a share of the optimal allocation's gain, in a single round with a fifth of the
population funded, as review noise increases. Filled circles: an equal division among
the top fifth of researchers by review score. Open circles: an equal division among
the top two fifths. Triangles: the partial lottery, which funds half of the top two
fifths at random. Dashed rule: uniform funding. Dotted rule: a lottery over all
researchers. In a heavy-tailed field with a tight budget (left), selection by review
gains most, and its gain declines slowly as review grows noisier; the partial lottery
gains less, mainly because it ignores the review ranking within the top two fifths;
the lottery over all researchers gains least, less than uniform funding. In an evenly
spread field with a large budget (right), the ordering reverses: uniform funding gains
most, and every scheme that concentrates grants gains less. Two thousand populations
per point.}
\label{fig:lottery-schemes}
\end{figure}


When does a lottery among the fundable lose little output? When the review ranking it
ignores is worth little. Replacing the review ranking within the top two fifths by
chance loses a few hundredths of the optimal allocation's gain in an evenly spread
field or when review is very noisy, and about a sixth of that gain in a heavy-tailed
field with sharp review.\footnote{The partial lottery's shortfall relative to
selection by review: 0.01 to 0.07 of the optimal allocation's gain in the evenly
spread field, and in the heavy-tailed field at the noisiest review; 0.14 to 0.18 in
the heavy-tailed field with sharp to moderately noisy review.} These are the fields in
which fine distinctions among applicants are worth little. The lottery debate rarely
asks a further question, whether to concentrate grants at all. In an evenly spread
field, our model's answer is to spread the budget thinly across many researchers, not
to hold a lottery among a fundable few. Whether a lottery's savings in review costs
and proposal effort outweigh the output it loses depends on both quantities. The
contest models cited above measure the savings; our model measures the loss.

Seed grants and lotteries both spread funds more widely than targeting does, and
behind both lies the older question of whether funding should be concentrated on a few
researchers or spread across many. Our results so far rest on one production form, in
which capability and resources are complements. One might hope that this choice
settles the concentration question: the harder it is to substitute resources for
capability, the more the money should go to the few researchers who can use it. To
test this hope we embed our production form in a family of technologies indexed by an
exponent $\gamma$, from Cobb-Douglas ($\gamma = 0$), under which a shortfall in either
input can be offset by more of the other, through our form ($\gamma = -1$), to Leontief,
under which output depends only on the scarcer input; Appendix~\ref{app:technology}
defines the family and reports which of our earlier results depend on the
choice.\footnote{The concentration runs are single-round, with two hundred seeds per
cell (four hundred at the refinement points of the evenly spread curve), and measure
the review-informed funder's grants, not the complete-information allocation. Leontief
allocations are computed with a finer allocation step, and ties at the kink are broken
by fill order.} The technology does not settle the question
(Figure~\ref{fig:concentration}). With a tight budget in a heavy-tailed field, harder
complementarity does concentrate the informed funder's grants. With an ample budget,
harder complementarity spreads the informed funder's grants, a pattern consistent with
a funder that raises every researcher toward their bottleneck once the budget allows
it. Where capability is evenly spread and the budget ample, the relation is not even
monotone: concentration peaks at an intermediate complementarity and declines toward
both Cobb-Douglas and Leontief.\footnote{Gini coefficients of the informed funder's
grants, from Cobb-Douglas to Leontief: 0.92 to 0.95 in a heavy-tailed field with a
tight budget ($\alpha_K = 1.3$, $b = 0.1$); 0.35 to 0.21 in the default field with an
ample budget ($b = 1$); 0.63 to 0.45 in a heavy-tailed field with an ample budget. In
the evenly spread field with an ample budget ($\alpha_K = 3.5$, $b = 1$) the Gini
coefficient is 0.16 at Cobb-Douglas, peaks at 0.196 near $\gamma = -6$, and falls to
0.185 at $\gamma = -12$; the peak's rise over the $\gamma = -12$ level is 3.5 standard
errors.} How widely funds should be spread is set jointly by the technology, the
budget, and the field's capability distribution; no one of the three decides it alone.

% Figure (S8): how concentrated the informed allocation is, across the family.
% y = Gini coefficient of the review-informed funder's grants (single round);
% x = CES exponent gamma (0 = Cobb-Douglas, -1 = our harmonic form, limit =
% Leontief, plotted as detached marks at the left edge). Three field-budget pairs.
% Ground: sigma_tierB_concentration (+ leontief endpoint, + 400-seed refinement for
% the evenly spread curve), 200 seeds per cell, re-derived in-session (verify_s9.R).
% Leontief cells carry the kink caveat (ties broken by fill order, n_steps = 800).
\begin{figure}[t]
\centering
\begin{tikzpicture}
\begin{axis}[
  width=0.82\textwidth, height=6.0cm,
  axis lines*=left,
  xmin=-14.8, xmax=0, ymin=-0.16, ymax=0.06,
  xtick={-14, -12, -6, -3, -1, 0},
  xticklabels={Leontief, $-12$, $-6$, $-3$, $-1$, $0$},
  ytick={-0.15, -0.10, -0.05, 0, 0.05},
  yticklabels={$-0.15$, $-0.10$, $-0.05$, 0, 0.05},
  xlabel={CES exponent $\gamma$ (complements $\leftarrow$, substitutable $\rightarrow$)},
  ylabel={change in grant concentration (Gini)},
  ylabel style={font=\small}, xlabel style={font=\small},
  tick label style={font=\small},
  axis line style={line width=0.4pt},
  tick style={line width=0.4pt, black},
  clip=false,
]
\draw[gray, dotted, line width=0.5pt] (axis cs:-1,-0.16) -- (axis cs:-1,0.06);
\node[font=\scriptsize, gray, anchor=south] at (axis cs:-1,0.061) {our form};
\addplot[black, line width=0.7pt, mark=*, mark size=1.1pt] table[x=g, y=dgini] {fig13-heavy01.dat};
\addplot[black, line width=0.7pt, mark=o, mark size=1.3pt] table[x=g, y=dgini] {fig13-def1.dat};
\addplot[black, line width=0.7pt, mark=triangle*, mark size=1.4pt] table[x=g, y=dgini] {fig13-even1.dat};
\addplot[black, only marks, mark=*, mark size=1.1pt] coordinates {(-14,0.031)};
\addplot[black, only marks, mark=o, mark size=1.3pt] coordinates {(-14,-0.147)};
\addplot[black, only marks, mark=triangle*, mark size=1.4pt] coordinates {(-14,0.009)};
\node[font=\scriptsize, anchor=north] at (axis cs:-8.5,0.014) {heavy-tailed, tight};
\node[font=\scriptsize, anchor=south] at (axis cs:-10.5,-0.047) {default, ample};
\node[font=\scriptsize, anchor=south] at (axis cs:-7.5,0.042) {evenly spread, ample};
\end{axis}
\end{tikzpicture}
\caption{Complementarity concentrates funding only when money is tight. The change in the Gini coefficient of the review-informed funder's grants, relative to its Cobb-Douglas level, as the production technology moves from Cobb-Douglas ($\gamma = 0$) toward perfect complements (Leontief, detached marks at the left edge). With a tight budget in a heavy-tailed field, stronger complementarity concentrates grants; with an ample budget it spreads them. Where capability is evenly spread and the budget ample, the relation is not even monotone: concentration peaks at an intermediate complementarity ($\gamma \approx -6$) and declines toward both ends. Cobb-Douglas levels: 0.92 (heavy-tailed, tight budget), 0.35 (default field, ample), 0.16 (evenly spread, ample). Two hundred simulated populations per point (four hundred at the refinement points of the evenly spread curve).}
\label{fig:concentration}
\end{figure}


Next we take stock: what the model implies for funders, and what its assumptions leave
out.

\begin{appendices}
% Appendix A: derivation of the gap rule (Proposition 1), with corollaries and
% examples. Source: for-claude/gap-rule-proof-v1.md (2026-08-14), converted 2026-09-08.
% Self-contained: states every assumption used; no KKT machinery (Lemma 3 does that
% work via the transfer argument and the concavity gradient inequality). Numerical
% verification: for-claude/verify_gap_rule.py (81 trials vs scipy SLSQP, machine
% precision; see the md's note 5).
\section{Derivation of the gap rule}\label{app:gap-rule}

\subsection*{The within-round allocation problem}

A population of $n$ researchers. Researcher $i$ has capability $K_i$ and baseline resources $R_{i0}$, both fixed within the round. A grant $g_i \geq 0$ adds to researcher $i$'s resources for the round, so that their expected output is $A\,\Lambda(K_i,\, R_{i0} + g_i)$ with $\Lambda(K, R) = 2KR/(K + R)$, the production function of \cref{sec:model}, and $A > 0$ the common productivity scale. Write
\[
\varphi_i(g) = \frac{2A\,K_i\,(R_{i0} + g)}{K_i + R_{i0} + g}
\]
for researcher $i$'s expected output as a function of their grant alone.

Throughout, $A > 0$, $K_i > 0$ and $R_{i0} \geq 0$ for every $i$, and the round's budget is $B > 0$. Nothing else about the model is used, so the result holds for every realized population and every round's budget; across rounds, the constant $c$ and the funded set are recomputed from the new population state.

The funder observes $(K_i, R_{i0})_{i=1}^n$ and divides the budget so as to maximize expected output:
\[
(\mathrm{P}) \qquad \max_{g_1, \ldots, g_n \geq 0} \; \sum_{i=1}^n \varphi_i(g_i) \quad \text{subject to} \quad \sum_{i=1}^n g_i \leq B.
\]

\subsection*{Three lemmas}

\begin{lemma}[the return to a grant]\label{lem:return}
For each $i$, on $g \in [0, \infty)$: (i) $\varphi_i$ is strictly increasing, with $\varphi_i'(g) = 2A\,K_i^2/(K_i + R_{i0} + g)^2 > 0$; (ii) $\varphi_i$ is strictly concave: $\varphi_i''(g) = -\,4A\,K_i^2/(K_i + R_{i0} + g)^3 < 0$; (iii) $\varphi_i'(g) < 2A$ whenever $R_{i0} + g > 0$, and $\varphi_i(g) \uparrow 2A K_i$ as $g \to \infty$.
\end{lemma}

\begin{proof}
Write $S = K_i + R_{i0} + g$. Then $\varphi_i(g) = 2AK_i(S - K_i)/S = 2AK_i - 2AK_i^2/S$, and (i) and (ii) follow by differentiating in $g$ (note $dS/dg = 1$). For (iii): $\varphi_i'(g) = 2A\,K_i^2/S^2 < 2A$ exactly when $S > K_i$, that is, when $R_{i0} + g > 0$; and $\varphi_i(g) = 2AK_i - 2AK_i^2/S \to 2AK_i$ as $S \to \infty$.
\end{proof}

\Cref{lem:return} says that the marginal value of a dollar to researcher $i$ increases with their capability and decreases with the resources they already hold, that each dollar granted lowers the value of the next, and that expected output never exceeds $2AK_i$.

\begin{lemma}[existence, uniqueness, exhaustion]\label{lem:existence}
Problem $(\mathrm{P})$ has exactly one solution $g^* = (g_1^*, \ldots, g_n^*)$, and it exhausts the budget: $\sum_i g_i^* = B$.
\end{lemma}
\begin{proof}
The feasible set is compact and convex and the objective $f(g) = \sum_i \varphi_i(g_i)$ is continuous and, by \cref{lem:return}(ii), strictly concave, so a maximizer exists and is unique. If $\sum_i g_i^* < B$, increasing any one $g_i^*$ slightly stays feasible and raises $f$ (\cref{lem:return}(i)), contradicting optimality.
\end{proof}
\begin{lemma}[equal marginal value]\label{lem:emv}
A feasible allocation $g$ with $\sum_i g_i = B$ solves $(\mathrm{P})$ if and only if there is a $\nu > 0$ such that
\[
\varphi_i'(g_i) = \nu \;\text{ for every } i \text{ with } g_i > 0, \qquad \varphi_i'(0) \leq \nu \;\text{ for every } i \text{ with } g_i = 0.
\]
\end{lemma}

\begin{proof}
\emph{Necessity.} Let $g^*$ solve $(\mathrm{P})$. Since $B > 0$ is exhausted (\cref{lem:existence}), some researcher is funded: $g_i^* > 0$ for some $i$. Fix any such $i$ and any $j \neq i$, and for $0 < \delta \leq g_i^*$ consider the transfer that replaces $g_i^*$ by $g_i^* - \delta$ and $g_j^*$ by $g_j^* + \delta$, leaving the rest unchanged. The result is feasible, so optimality gives
\[
\varphi_j(g_j^* + \delta) - \varphi_j(g_j^*) \;\leq\; \varphi_i(g_i^*) - \varphi_i(g_i^* - \delta).
\]
Dividing by $\delta$ and letting $\delta \downarrow 0$ yields $\varphi_j'(g_j^*) \leq \varphi_i'(g_i^*)$. If $j$ is also funded, the same argument with $i$ and $j$ exchanged gives the reverse inequality, so all funded researchers share a common marginal value; call it $\nu = \varphi_i'(g_i^*) > 0$ (positivity by \cref{lem:return}(i)). For unfunded $j$, the displayed inequality reads $\varphi_j'(0) \leq \nu$.

\emph{Sufficiency.} Let $g$ satisfy the condition with multiplier $\nu$, and let $h$ be any feasible allocation. Concavity of $\varphi_i$ gives, for every $i$, $\varphi_i(h_i) \leq \varphi_i(g_i) + \varphi_i'(g_i)\,(h_i - g_i)$. For funded $i$, $\varphi_i'(g_i) = \nu$, so $\varphi_i'(g_i)(h_i - g_i) = \nu\,(h_i - g_i)$. For unfunded $i$, $g_i = 0$ and $h_i - g_i = h_i \geq 0$, so $\varphi_i'(0)\,(h_i - g_i) \leq \nu\,(h_i - g_i)$. Summing over $i$:
\[
f(h) \;\leq\; f(g) + \nu \sum_i (h_i - g_i) \;=\; f(g) + \nu\Big(\sum_i h_i - B\Big) \;\leq\; f(g),
\]
since $h$ is feasible and $\nu > 0$. So $g$ solves $(\mathrm{P})$.
\end{proof}

\Cref{lem:emv} is the formal content of the sentence in \cref{sec:optimal}: an allocation is optimal exactly when no dollar can be moved to a researcher who can produce more value with it.

\subsection*{The proposition}

\begin{proprestate}[the gap rule]
The unique solution of $(\mathrm{P})$ is
\[
g_i^* = \max(c\,K_i - R_{i0},\, 0), \qquad i = 1, \ldots, n,
\]
where $c > 0$ is the unique constant satisfying the budget identity $\sum_{i=1}^n \max(c\,K_i - R_{i0},\, 0) = B$. Equivalently, $c = \sqrt{2A/\nu} - 1$, where $\nu \in (0, 2A)$ is the common marginal value of \cref{lem:emv}.
\end{proprestate}

\begin{proof}
Let $g^*$ be the unique solution (\cref{lem:existence}) and $\nu$ the multiplier that \cref{lem:emv} associates with it. The proof proceeds in four steps.

\emph{Step 1: $\nu \in (0, 2A)$, so $c = \sqrt{2A/\nu} - 1$ is well defined and positive.} Positivity of $\nu$ is part of \cref{lem:emv}. For the upper bound: some researcher $i$ is funded, and for them $R_{i0} + g_i^* \geq g_i^* > 0$, so \cref{lem:return}(iii) gives $\nu = \varphi_i'(g_i^*) < 2A$. Then $2A/\nu > 1$, so $c = \sqrt{2A/\nu} - 1 > 0$.

\emph{Step 2: funded researchers receive exactly their gap.} Let $g_i^* > 0$. By \cref{lem:emv}, $2A\,K_i^2/(K_i + R_{i0} + g_i^*)^2 = \nu$. Solving, and taking the positive square root (both sides of the next display are positive),
\[
K_i + R_{i0} + g_i^* = K_i\,\sqrt{2A/\nu} = (1 + c)\,K_i,
\]
so $g_i^* = c\,K_i - R_{i0}$. Since $g_i^* > 0$, this equals $\max(c\,K_i - R_{i0},\, 0)$.

\emph{Step 3: unfunded researchers have no gap.} Let $g_i^* = 0$. By \cref{lem:emv}, $2A\,K_i^2/(K_i + R_{i0})^2 = \varphi_i'(0) \leq \nu$, so $K_i + R_{i0} \geq K_i\sqrt{2A/\nu} = (1 + c)\,K_i$, that is, $R_{i0} \geq c\,K_i$. Then $\max(c\,K_i - R_{i0},\, 0) = 0 = g_i^*$.

Steps 2 and 3 establish $g_i^* = \max(c\,K_i - R_{i0}, 0)$ for every $i$, and the budget identity is exhaustion (\cref{lem:existence}). It remains to show $c$ is the only constant satisfying that identity.

\emph{Step 4: the budget identity pins $c$ uniquely.} Define $G(c) = \sum_{i=1}^n \max(c\,K_i - R_{i0},\, 0)$ for $c \geq 0$. $G$ is continuous and nondecreasing, with $G(0) = 0$ and $G(c) \to \infty$ as $c \to \infty$. Moreover $G$ is strictly increasing wherever it is positive: if $G(c) > 0$, then $c\,K_j > R_{j0}$ for some $j$, and for any $c' > c$, $G(c') \geq G(c) + K_j\,(c' - c) > G(c)$. Now suppose $c_1 < c_2$ both satisfy $G(c) = B$. Since $B > 0$, $G(c_1) > 0$, so $G(c_2) > G(c_1) = B$, a contradiction.
\end{proof}

\subsection*{Consequences}

\begin{corollary}[funding frontier and budget comparative statics]\label{cor:frontier}
(i) Researcher $i$ is funded, $g_i^* > 0$, exactly when $K_i > R_{i0}/c$. (ii) As a function of the budget $B$, the constant $c(B)$ is continuous and strictly increasing; hence the funded set is nondecreasing in $B$, every funded researcher's grant $c(B)\,K_i - R_{i0}$ is strictly increasing in $B$, and the common marginal value $\nu(B) = 2A/(1 + c(B))^2$ is strictly decreasing in $B$. Both parts follow from the formula and Step 4.
\end{corollary}
\begin{corollary}[the value of targeting vanishes in the budget]\label{cor:targeting-vanishes}
Let $f^*(B)$ denote the optimal value of $(\mathrm{P})$ and $f^u(B) = \sum_i \varphi_i(B/n)$ the value of uniform funding. Then $0 \leq f^*(B) - f^u(B)$ for every $B$, and $f^*(B) - f^u(B) \to 0$ as $B \to \infty$.
\end{corollary}

\begin{proof}
The first inequality holds because uniform funding is feasible. For the second: by \cref{lem:return}(iii), $\varphi_i(g) < 2AK_i$ for every $g$, so $f^*(B) \leq 2A\sum_i K_i$; and $f^u(B) = \sum_i \varphi_i(B/n) \to 2A\sum_i K_i$ as $B \to \infty$. Hence $0 \leq f^*(B) - f^u(B) \leq 2A\sum_i K_i - f^u(B) \to 0$.
\end{proof}

Since uniform funding's own gain over no funding does not vanish (it approaches $2A\sum_i K_i - \sum_i \varphi_i(0) > 0$), the same conclusion holds for the value of targeting measured relative to that gain, the quantity plotted in Figure~\ref{fig:targeting-value}.

\begin{example}[funding the track record can produce less than uniform funding]\label{ex:track-record}
Two researchers, $A = 1$, budget $B = 2$. Researcher 1: $K_1 = 8$, $R_{10} = 8$ (productive and well resourced; expected output $\varphi_1(0) = 8$). Researcher 2: $K_2 = 4$, $R_{20} = 1$ (capable but resource-starved; $\varphi_2(0) = 8/5$). Allocating in proportion to expected output gives $g = (5/3,\, 1/3)$ and total output $570/53 \approx 10.75$. Uniform funding, $g = (1, 1)$, gives $568/51 \approx 11.14$. The gap rule ($c = 3/4$) gives $g^* = (0, 2)$ and $80/7 \approx 11.43$.
\end{example}

\begin{example}[funding the under-resourced can produce less than uniform funding]\label{ex:scarcity}
Two researchers, $A = 1$, budget $B = 2$. Researcher 1: $K_1 = 10$, $R_{10} = 5$. Researcher 2: $K_2 = 1/2$, $R_{20} = 0$ (the scarcest resources in the population). Directing the budget to the most resource-starved researcher gives total output $112/15 \approx 7.47$. Uniform funding gives $49/6 \approx 8.17$. The gap rule ($c = 2/3$) gives $g^* = (5/3,\, 1/3)$ and $42/5 = 8.4$.
\end{example}

% Appendix B: lotteries are outperformed by equal division (the Jensen argument).
% Referenced by S8; \input inside the appendices environment AFTER appendix-gap-rule
% when S8 is typeset. Uses the prop counter (renders as Proposition 2) and Lemma 1
% of Appendix A.
\section{Lotteries and equal division}\label{app:lottery}

Fix a round, a budget $B > 0$, and a pool $P \subseteq \{1, \ldots, n\}$ of $m$
researchers. The \emph{equal division} over $P$ grants each researcher in $P$ the
amount $B/m$. A \emph{lottery} over $P$ with $k \leq m$ winners selects $k$
researchers from $P$ uniformly at random and grants each winner $B/k$. Expected
output is evaluated as in Appendix~\ref{app:gap-rule}: given an allocation $g$, it is
$\sum_i \varphi_i(g_i)$, and for a random allocation, the expectation is taken over
the draw as well.

\begin{prop}[lotteries are outperformed by equal division]\label{prop:lottery}
For every pool $P$, budget $B$, and winner count $k \leq m$, the lottery's expected
output does not exceed the equal division's, with equality only if $k = m$. In
particular, a lottery over the whole population, at any winner count, produces at
most uniform funding's expected output.
\end{prop}

\begin{proof}
Researchers outside $P$ receive nothing under either scheme. For $i \in P$, the
lottery's grant to $i$ is a random variable $\tilde g_i$ equal to $B/k$ with
probability $k/m$ and $0$ otherwise, so $\mathbb{E}[\tilde g_i] = B/m$. By linearity,
the lottery's expected output is $\sum_{i \notin P} \varphi_i(0) + \sum_{i \in P}
\mathbb{E}[\varphi_i(\tilde g_i)]$. Since $\varphi_i$ is strictly concave
(Lemma~\ref{lem:return}(ii)), Jensen's inequality gives
$\mathbb{E}[\varphi_i(\tilde g_i)] \leq \varphi_i(\mathbb{E}[\tilde g_i]) =
\varphi_i(B/m)$ for every $i \in P$, with equality only if $\tilde g_i$ is
degenerate, that is, $k = m$. Summing over $i$ yields the claim; taking
$P = \{1, \ldots, n\}$ yields the comparison with uniform funding.
\end{proof}

\begin{remark}[what the proposition does and does not cover]
The proposition compares expected output within a fixed pool. It measures the effect
of the randomization alone: whatever pool review selects, replacing the lottery within
that pool by the equal division of the same money weakly increases expected output.
The proposition says nothing about the choice of pool, about the review ranking
within the pool that the equal division also ignores, or about anything a lottery
saves outside the model (review costs, proposal effort, bias). \Cref{sec:egalitarian}
measures the first two and cites the literature on the third.
\end{remark}

% Appendix C: robustness to the production technology.
% Referenced by S3, S6, S7, S8. \input inside the appendices environment AFTER
% appendix-lottery. Holds Figure 12 (signal robustness) and Table 2 (timing by
% technology). Grounds: sigma_tierA_gc0 (hash d77c854, 100 seeds), sigma_tierA_gcm3
% (100 seeds), sigma_tierA_leontief (50 seeds), all re-derived in-session
% (verify_s9.R); the tier A sweeps use an earlier version of the allocation routine
% than the body's runs, so magnitudes are compared within this appendix only.
\section{Robustness to the production technology}\label{app:technology}

Throughout the paper, a researcher's expected output has taken the form of \S\ref{sec:model}, $2AKR/(K + R)$, in which capability and resources are complements and output is limited by the scarcer of the two. This appendix asks which of our results depend on that choice. We embed our form in a family of production technologies indexed by an exponent $\gamma$: expected output is $A\,M_\gamma(K, R)$, where $M_\gamma(K, R) = \big((K^{\gamma} + R^{\gamma})/2\big)^{1/\gamma}$ is the power mean of capability and resources. At $\gamma = -1$ the power mean is the harmonic mean $2KR/(K + R)$, and the family reproduces our form. In the limit $\gamma \to 0$, the Cobb-Douglas case, the power mean is the geometric mean $\sqrt{KR}$, and a shortfall in either input can be offset by more of the other; as $\gamma$ decreases the inputs grow harder to substitute for one another; and in the limit $\gamma \to -\infty$, the Leontief case, the power mean is $\min(K, R)$, and output depends only on the scarcer input. We re-run the model's central comparisons across this family.\footnote{Two rounds, one hundred simulated populations per point (fifty for Leontief). Leontief allocations are computed with a finer allocation step, and ties at the kink are broken by fill order; the Leontief points carry that caveat. These runs use an earlier version of the allocation routine than the runs reported in the body, so their magnitudes are compared with one another and not with the body's figures.}

The value of the review signal survives the whole family (Figure~\ref{fig:signal-robustness}). At every technology tested, from Cobb-Douglas to Leontief, the signal's value increases as capability grows more heavy-tailed and as review grows more informative. The ordering of review's value by field (\S\ref{sec:review}) is therefore not an artifact of our choice of production function. What changes with the technology is the scale. The more complementary the inputs, the more review is worth: under heavy-tailed capability with sharp review, the signal's value increases from about seven percent of no-funding output at Cobb-Douglas to about thirty percent at $\gamma = -3$ and a third at Leontief.\footnote{6.8, 29.3, and 33.7 percent of no-funding output; the same ordering holds at every tested tail and noise level at which the value is distinguishable from zero.} Complementarity raises what targeting is worth, and with it the value of the information that targeting requires.

% Figure (Appendix C): the structure of review's value survives the production family.
% y = value of the review signal (S5 - S4) as percent of no-funding output.
% Left panel: vs the capability tail parameter alpha_K at sharp review (tau = 0.3).
% Right panel: vs review noise tau at heavy tails (alpha_K = 1.3).
% Curves: Cobb-Douglas (gamma = 0, filled circles), gamma = -3 (open circles),
% Leontief (triangles). Ground: sigma_tierA_{gc0,gcm3,leontief} signal_value
% summaries (T = 2, greedy allocator n_steps = 400; 100/100/50 seeds), re-derived
% in-session (verify_s9.R). Data: fig12-left.dat, fig12-right.dat.
\begin{figure}[t]
\centering
\begin{tikzpicture}
\begin{axis}[
  name=left,
  width=0.46\textwidth, height=5.6cm,
  axis lines*=left,
  xmin=1.3, xmax=3.5, ymin=0, ymax=35,
  xtick={1.3, 1.8, 2.5, 3.5},
  ytick={0, 10, 20, 30},
  xlabel={capability tail $\alpha_K$},
  ylabel={value of review (\% of no-funding output)},
  ylabel style={font=\small}, xlabel style={font=\small},
  tick label style={font=\small},
  axis line style={line width=0.4pt},
  tick style={line width=0.4pt, black},
  title={\small sharp review ($\tau_K = 0.3$)},
  title style={yshift=-2pt},
  clip=false,
]
\addplot[black, line width=0.7pt, mark=*, mark size=1.1pt] table[x=k, y=cd] {fig12-left.dat};
\addplot[black, line width=0.7pt, mark=o, mark size=1.3pt] table[x=k, y=g3] {fig12-left.dat};
\addplot[black, line width=0.7pt, mark=triangle*, mark size=1.4pt] table[x=k, y=leo] {fig12-left.dat};
\node[font=\scriptsize, anchor=south west] at (axis cs:1.33,7.3) {Cobb-Douglas};
\end{axis}
\begin{axis}[
  at={(left.outer east)}, anchor=outer west, xshift=6pt,
  width=0.46\textwidth, height=5.6cm,
  axis lines*=left,
  xmode=log,
  xmin=0.3, xmax=10, ymin=0, ymax=35,
  xtick={0.3, 1, 3, 10},
  xticklabels={0.3, 1, 3, 10},
  ytick={0, 10, 20, 30},
  yticklabels={,,,},
  xlabel={review noise $\tau_K$},
  xlabel style={font=\small},
  tick label style={font=\small},
  axis line style={line width=0.4pt},
  tick style={line width=0.4pt, black},
  title={\small heavy-tailed ($\alpha_K = 1.3$)},
  title style={yshift=-2pt},
  clip=false,
]
\addplot[black, line width=0.7pt, mark=*, mark size=1.1pt] table[x=tau, y=cd] {fig12-right.dat};
\addplot[black, line width=0.7pt, mark=o, mark size=1.3pt] table[x=tau, y=g3] {fig12-right.dat};
\addplot[black, line width=0.7pt, mark=triangle*, mark size=1.4pt] table[x=tau, y=leo] {fig12-right.dat};
\node[font=\scriptsize, anchor=south west] at (axis cs:0.315,7.3) {Cobb-Douglas};
\end{axis}
\end{tikzpicture}
\caption{The structure of review's value survives the production technology. The value of the review signal as a percent of no-funding output, across the family of production technologies (the CES family): at every technology tested, the value increases as capability grows more heavy-tailed (left, decreasing $\alpha_K$) and as review grows more informative (right, decreasing $\tau_K$), and its scale increases with complementarity. Filled circles: Cobb-Douglas; open circles: $\gamma = -3$; triangles: Leontief. Two rounds; one hundred simulated populations per point (fifty for Leontief).}
\label{fig:signal-robustness}
\end{figure}


The timing results of \S\ref{sec:timing} do not survive the substitutable end of the family (Table~\ref{tab:timing-technology}). Under Cobb-Douglas production the value of scheduling is indistinguishable from zero at every compounding rate, and the optimal schedule stays within about a hundredth of even. At $\gamma = -3$, a technology more complementary than our form, both the value of scheduling and the lean toward late spending increase with the compounding rate and exceed their Cobb-Douglas values at every rate. The timing results hold where capability and resources are complements and vanish where they are substitutable.

\begin{table}[t]
\centering
\small
\caption{The value of scheduling and the optimal schedule by production technology. Five rounds. The value of scheduling is the forward-looking funder's gain over the myopic funder, both with review, as a percent of no-funding output; the schedule's center of mass is 0.5 when spending is even and larger when spending leans late. One hundred simulated populations per point; standard errors of the gain are at most 0.002 (Cobb-Douglas) and 0.09 ($\gamma = -3$) percentage points.}
\label{tab:timing-technology}
\begin{tabular}{@{}lccccc@{}}
\toprule
compounding rate $\epsilon$ & 0.05 & 0.15 & 0.30 & 0.55 & 0.85 \\
\midrule
\multicolumn{6}{@{}l}{\emph{Value of scheduling (percent of no-funding output)}} \\
Cobb-Douglas ($\gamma = 0$) & $-0.001$ & $-0.002$ & $0.001$ & $0.001$ & $0.004$ \\
$\gamma = -3$               & $0.01$   & $0.13$   & $0.45$  & $1.07$  & $1.54$ \\
\addlinespace
\multicolumn{6}{@{}l}{\emph{Center of mass of the optimal schedule}} \\
Cobb-Douglas ($\gamma = 0$) & 0.499 & 0.500 & 0.501 & 0.505 & 0.512 \\
$\gamma = -3$               & 0.511 & 0.540 & 0.577 & 0.623 & 0.660 \\
\bottomrule
\end{tabular}
\end{table}

How widely the informed funder spreads its grants across this family is reported in \S\ref{sec:egalitarian} (Figure~\ref{fig:concentration}).

\end{appendices}

\bibliographystyle{chicago}
\bibliography{references}

\end{document}
